\documentclass[aps,prd,twocolumn,showpacs,superscriptaddress,preprintnumbers]{revtex4-2}
\usepackage{amsmath,amssymb,amsfonts,amsthm}
\usepackage{graphicx}
\usepackage{hyperref}
\usepackage{dsfont}
\usepackage{bm}
\usepackage{booktabs}

\newtheorem{lemma}{Lemma}[section]
\newtheorem{theorem}[lemma]{Theorem}
\renewcommand{\proofname}{Proof}

\begin{document}

\title{An Information-Capacity Model of Gravity and Quantum Decoherence}

\author{Huang Zhongchang}
\affiliation{Independent Researcher, Nanjing, China}

\date{\today}

\begin{abstract}
We construct a model in which Newtonian gravity and quantum decoherence are modeled as regimes of a single dynamical equation. The model has one irreducible definition---a causal event is an asymmetric influence producing a determinate change in a discrete cell (Definition~C)---plus a small set of motivated structural choices (regular lattice, logarithmic gradient, linear damping ansatz, universality class $\Phi$; enumerated in Table~I). From these ingredients we obtain the unified telegraph equation $\partial_t^2 q + \gamma_0(1-q)\partial_t q = c^2\nabla^2(\ln q)$ for the mesoscopic free-capacity fraction $q$. The static limit $\nabla^2(\ln 1/q)=0$ recovers the $1/r$ functional form of the Newtonian potential; the coefficient is calibrated to the observed value of $G$. In vacuum ($q\to1$), state-dependent damping $\gamma(q)=\gamma_0(1-q)$ vanishes, restoring the Lorentz-invariant sector that the model always contained; preferred-frame effects are confined to matter-filled regions and suppressed by $GM/rc^2\sim10^{-6}$ at the solar surface. We prove an H-theorem via a Lyapunov functional, and we show that the static equation decomposes exactly into a diffusion term ($\nabla^2 q/q$) and a self-organization term ($|\nabla q|^2/q^2$), unifying decoherence and gravitational collapse as two regimes of one process.

We enumerate nine limitations: (1)~the damping form is a modeling ansatz, not uniquely derived; (2)~$G$ is calibrated, not predicted; (3)~the spacetime metric is a structural postulate (causal locking hypothesis), not derived from Definition~C; (4)~only the Newtonian limit of GR is recovered; (5)~matter coupling is not microscopically derived; (6)~the 1-bit quantization argument is structural, not a formal theorem; (7)~GW170817 consistency requires an unproven tensor-scalar decoupling conjecture; (8)~the volume-scaling entropy bound conflicts with the holographic area-scaling bound; (9)~the continuum limit and $d=3$ are empirical inputs.

The paper's contribution is not a unique derivation of gravity from causal structure, but a demonstration of a model---compatible with Definition~C---in which gravity-like and decoherence-like phenomena coexist within a single equation, with a well-defined inventory of what is derived, what is chosen, and what remains open.
\end{abstract}

\maketitle

%%%%%%%%%%%%%%%%%%%%%%%%%%%%%%%%%%%%%%%%%%%%%%%%%%%%%%%%%%%%%%%%%%%%%%
\section{Introduction}
\label{sec:intro}
%%%%%%%%%%%%%%%%%%%%%%%%%%%%%%%%%%%%%%%%%%%%%%%%%%%%%%%%%%%%%%%%%%%%%%

Unifying quantum mechanics with gravity. Few problems in fundamental physics have resisted solution for so long.
Quantum field theory---built by quantizing classical fields on a fixed spacetime background---is the most precisely tested framework in physics, yet general relativity, describing gravity as spacetime curvature, has passed every observational test~\cite{Will:2014}.
The two frameworks are structurally incompatible: GR is a classical, deterministic, geometric theory; quantum mechanics is a probabilistic theory on Hilbert space.
String theory~\cite{Green:1987sp,Polchinski:1998rq}, loop quantum gravity~\cite{Rovelli:2004tv,Ashtekar:2004eh}, causal dynamical triangulations~\cite{Ambjorn:2001cv}, asymptotic safety~\cite{Weinberg:1979,Reuter:2012}, entropic gravity~\cite{Verlinde:2010hp,Jacobson:1995ab}, and the quantum memory matrix~\cite{Neukart:2024qmm}---numerous approaches exist.
None has yielded a complete, experimentally verified theory.

This paper takes a different approach. Rather than quantizing geometry or geometrizing quantum fields, we ask a simpler question: can gravity be derived from \emph{information-theoretic} first principles?
Put differently, if both quantum mechanics and gravity are consequences of the same underlying constraint---the finite information capacity of spacetime---then their apparent incompatibility may dissolve.
That is the hypothesis we explore here.

Our quantum-mechanical foundation rests on a body of peer-reviewed work establishing that finite information capacity plus structural constraints yields standard quantum mechanics. Zeilinger (1999)~\cite{Zeilinger1999} identified that an elementary system carries exactly one bit of information, from which irreducible randomness and entanglement follow. Daki\'c and Brukner (2011)~\cite{Dakic2011} formalized the 1-bit capacity as an axiom and proved quantum theory is the unique non-classical probabilistic theory satisfying it. Chiribella, D'Ariano, and Perinotti (2011)~\cite{Chiribella2011} derived the full Hilbert-space structure and Born rule from information-theoretic axioms without assuming quantum primitives. Masanes and M\"uller (2011)~\cite{Masanes2011} proved that the purification postulate singles out quantum theory uniquely among general probabilistic theories. Shaghaghi (2023)~\cite{Shaghaghi2023} showed that single-variable systems carry at most one bit, from which the Born rule and Schr\"odinger equation follow. Zilly (2026)~\cite{Zilly:2026} independently reached the same conclusions via graded equality, providing a machine-checked proof in Lean~4; this result, while recent and awaiting community verification, is consistent with the established reconstruction literature.

We accept this body of work as the quantum-mechanical foundation. It means quantum mechanics is itself a consequence of bounded capacity, not a separate axiom. $\hbar$ enters as the scale of the capacity bound, fixing phase space volume quantization~\cite{Zilly:2026,Shaghaghi2023}.
We build upward from this foundation toward gravity---a construction that takes QM as the already-derived starting point and asks what else falls out of the same scaffolding.

\textbf{Quantum foundations: what we rely on.}
The identification of the cell spacing $a$ with the Planck length $\ell_P$ uses $\hbar$---an empirical constant (Tier~1) whose role in phase space quantization is independently established by the reconstruction literature~\cite{Zeilinger1999,Dakic2011,Chiribella2011,Shaghaghi2023}. The physical interpretation of the 1-bit capacity bound as connected to quantum coherence is supported by the same body of work. The classical $q$-field equations are self-contained; their quantum interpretation is not, but it rests on peer-reviewed foundations rather than a single preprint.
Zilly's theorem provides the quantum-mechanical context but is not load-bearing for the derivation of the unified telegraph equation, vacuum Lorentz restoration, or Newtonian gravity.

We introduce a \textbf{single definition}:

\begin{quote}
\textbf{Definition~C (Causal Event).} A causal event is an asymmetric influence that produces a determinate change in a discrete physical cell.
\end{quote}

Asymmetry provides the arrow of time. Determinate change---formalized as a single-valued transition map $T_e: s \mapsto T_e(s)$ on the cell's state space---is shown, via a structural argument (Sec.~\ref{sec:axioms}.A), to imply that every elementary causal event transfers exactly 1 bit. The 1-bit quantization is not an additional postulate; it follows from requiring that the transition map be single-valued (determinate), non-identity (something changes), and elementary (not decomposable). Finite propagation at speed $c$ entails locality. Between causal events, a cell's state is objectively undetermined---this undeterminedness \emph{is} quantum superposition, whose mathematical form is provided by Zilly's theorem.

Section~\ref{sec:axioms}.A provides the full 1-bit structural argument, inventories its vulnerabilities, and explains the division of labor between Definition~C and the QM reconstruction literature.

From this definition we construct a dynamical theory of the free-capacity fraction $q_i(t)$, the fraction of cell $i$'s 1-bit capacity that remains available for new information.

\textbf{QM independence of the classical sector.} The quantum-mechanical reconstruction literature~\cite{Zeilinger1999,Dakic2011,Chiribella2011,Masanes2011,Shaghaghi2023,Zilly:2026} assures us that standard QM is compatible with---and derivable from---finite-capacity principles. However, \emph{no result in Sections~\ref{sec:dynamics}--\ref{sec:selforg} depends on these QM foundations.} The derivation of the unified telegraph equation, the dynamical Lorentz restoration, the static Newtonian limit, and the H-theorem require only Definition~C and the Tier~2 structural choices enumerated in Table~\ref{tab:tiers}. The QM foundations enter solely to connect the classical gravitational limit back to the quantum-mechanical interpretation, closing the conceptual loop---but the classical derivation stands independently.

The main results are:

\begin{itemize}
    \item The unified current equation with $\bar{q}=(q_i+q_j)/2$ and the continuity equation combine to yield $\partial_t^2 q + \gamma_0(1-q)\partial_t q = c^2\nabla^2(\ln q)$. For small perturbations around $q=1$, $\ln q\approx q-1$ recovers the linear telegraph equation; the $\gamma\to0$ vacuum limit yields $\Box(\ln q)=0$; and the static limit directly gives $\nabla^2(\ln q)=0$.
    \item In vacuum ($q\to1$), $\gamma\to0$ and Lorentz invariance is restored in vacuum by the vanishing of state-dependent damping, without being separately imposed as an axiom.
    \item The preferred-frame effects from $\gamma>0$ are confined to matter-filled regions ($q<1$), where they are suppressed by $(1-q)\sim GM/rc^2\sim10^{-6}$ near the solar surface---consistent with all Lorentz-violation constraints.
    \item In the static limit, $\nabla^2(\ln 1/q)=0$, whose spherically symmetric solution $q(r)=e^{-GM/rc^2}$ recovers the Newtonian potential $\Phi=-GM/r$ at large $r$.
    \item The decomposition $\nabla^2 q/q = |\nabla q|^2/q^2$ unifies quantum decoherence (diffusion-dominant at $q\approx1$) and gravitational collapse (self-organization-dominant at $q\approx0$) as two regimes of a single information-dynamic process.
    \item An H-theorem for the full nonlinear theory is proved via a Lyapunov functional, establishing that the dynamics are thermodynamically consistent.
\end{itemize}

This work builds on and differs from prior approaches in specific ways.
Jacobson~\cite{Jacobson:1995ab} derived the Einstein equations from thermodynamic considerations applied to local Rindler horizons; we go further, deriving the Newtonian potential from information dynamics without assuming any metric structure.
Verlinde~\cite{Verlinde:2010hp} proposed entropic gravity; we provide the microscopic mechanism---the $q$-field and $J$-current---that his approach lacks.
Padmanabhan~\cite{Padmanabhan:2010} explored the thermodynamic structure of gravitational field equations; our framework reproduces similar structure from a single information-theoretic definition and proves an explicit H-theorem.
Bassani and Magueijo~\cite{Bassani:2025htma} showed that diffeomorphism invariance can emerge from a Markov chain of stochastic choices; our Lorentz selection is deterministic and requires far fewer assumptions, though the emergent symmetry is weaker (Lorentz vs.\ diffeomorphism).
We provide a detailed, methodologically fair comparison in Sec.~\ref{sec:discussion}, applying the same tiered-assumption methodology to both theories.

The paper is organized as follows.
Section~\ref{sec:axioms} presents the causal definition and constructs the cell graph, including a tiered classification of all assumptions.
Section~\ref{sec:dynamics} derives the unified $q$-field dynamics from the fundamental current equation.
Section~\ref{sec:gamma} constructs $\gamma(q)=\gamma_0(1-q)$ from the causal definition and frames $\gamma_0$ through naturalness.
Section~\ref{sec:lorentz} proves that Lorentz invariance is restored in vacuum.
Section~\ref{sec:static} takes the static limit and recovers Newtonian gravity through the unified derivation.
Section~\ref{sec:selforg} analyzes self-organization, proves the H-theorem, and gives the stability analysis.
Section~\ref{sec:predictions} presents testable predictions.
Section~\ref{sec:discussion} discusses limitations, compares with HTMAU using tiered methodology, and outlines open directions.

%%%%%%%%%%%%%%%%%%%%%%%%%%%%%%%%%%%%%%%%%%%%%%%%%%%%%%%%%%%%%%%%%%%%%%
\section{Causal Definition and Cell Graph}
\label{sec:axioms}
%%%%%%%%%%%%%%%%%%%%%%%%%%%%%%%%%%%%%%%%%%%%%%%%%%%%%%%%%%%%%%%%%%%%%%

\subsection{The Causal Definition}
\label{sec:causal_def}

We begin with a single definition.

\begin{quote}
\textbf{Definition~C (Causal Event).} A causal event is an asymmetric influence that produces a determinate change in a discrete physical cell.
\end{quote}

This definition contains only two primitive concepts---\emph{asymmetry} and \emph{determinate change}---plus the discrete cell as the locus of the event. Much follows from it. The remainder of this subsection unpacks the definition and shows that the 1-bit quantization, locality, and quantum superposition are not additional assumptions layered on top of it; they are consequences of what the definition asserts.

\textbf{Asymmetry $\Rightarrow$ time arrow.}
If A influences B in a causal event, B does not simultaneously influence A in that same event. This asymmetry is the origin of temporal directionality: the directed edge from cause to effect defines a partial order, and the accumulation of such edges constitutes what we call time. Without asymmetry, there is no distinction between cause and effect, and the word ``causal'' has no content. The arrow of time is not an extra axiom; it is the direction of the arrow inscribed in each causal event.

\textbf{Determinate change $\Rightarrow$ 1-bit quantization.}
This is the central deduction. We present it as a structural argument---not a formal mathematical theorem---because the link from ``determinate'' to ``discrete'' relies on a physical premise (finite distinguishability at the Planck scale) that is well-motivated but not logically forced by Definition~C alone. The argument proceeds in four steps.

Fix a target cell $B$. Let $\Sigma$ be its state space. A causal event $e: A \to B$ defines a transition map
\begin{equation}
T_e: \Sigma \to \Sigma, \qquad s \mapsto T_e(s),
\label{eq:transition}
\end{equation}
where $s$ is the state of $B$ before the event and $T_e(s)$ is its state after. For $e$ to qualify as a causal event, $T_e$ must be \textbf{single-valued} (a function, not a relation) and \textbf{non-identity} ($T_e \neq \mathrm{id}$): the event produces a determinate change.

\begin{enumerate}
    \item \textbf{Determinate change is discrete change.} Definition~C specifies a \emph{discrete} physical cell. The word ``discrete'' is part of the definition: each cell is a countable entity among finitely or countably many such entities in any finite spatial region. If a cell had a continuum of distinguishable internal states, specifying which state obtains would require infinite precision---infinite information. No single causal event can convey infinite information; if it could, it would be indistinguishable from infinitely many events. A single-valued transition on a discrete cell therefore maps to one of finitely many distinguishable outcomes. A continuum objection (``a Gaussian measure gives finite Shannon entropy for a continuum'') is inapplicable: the causal event $e$ itself provides no probability measure on $\Sigma$; the measure would be an external input beyond Definition~C.

    \item \textbf{A finite set of outcomes has a minimal non-identity change.} Since $O = \{T_e(s) : s \in \Sigma\}$ is finite (step 1), the set of change magnitudes $\Delta(s) = \text{\# of distinguishable state distinctions between } s \text{ and } T_e(s)$ is a finite set of non-negative integers. Excluding zero (no change), a finite set of positive integers attains a minimum $m$.

    \item \textbf{The minimum is exactly 2 states.} A change distinguishing only 1 state is no change at all (before and after are the same). The smallest possible distinguishable change therefore distinguishes exactly 2 states: the pre-event state and its distinct image. Two states define one bit of information: the capacity of a binary channel is $\log_2 2 = 1$ bit~\cite{Shannon1948}, independent of outcome probabilities.

    \item \textbf{Every elementary causal event is a 1-bit-capacity event.} A causal event with $|O| = 2^n$ ($n > 1$) distinguishable outcomes is---in the classical setting where bits are independently identifiable---a composite of $n$ elementary 1-bit events. What about $|O| \neq 2^n$ (e.g., 3, 5, 6)? If the outcome set has a non-power-of-2 cardinality, it cannot be encoded in an integer number of bits without redundancy. But physical distinguishability is binary at root (a state either is or is not the given outcome), and any finite set of $k$ distinguishable alternatives can be encoded in $\lceil\log_2 k\rceil$ bits. The minimal change ($|O|=2$, $\lceil\log_2 2\rceil = 1$) remains the irreducible unit; larger outcome sets are compounds with redundant encodings. The quantum case (irreducibly entangled alternatives) is the recognized exception; Zilly's theorem covers it without increasing per-cell capacity.
\end{enumerate}

\textbf{Status of this argument.} The deduction is a structural argument, not a formal theorem. It is vulnerable at three points:
\begin{enumerate}
    \item \textbf{Single-valuedness.} An inspector may propose relational (indeterminate) causality. We reject this because ``$A$ caused $B$ to change'' without a specific ``to what'' lacks a truth value. But the rejection is methodological, not logically forced.
    \item \textbf{Discrete $\Rightarrow$ finite states.} A discrete set can be infinite (e.g., $\mathbb{N}$). The step from ``discrete cell'' to ``finitely many distinguishable outcomes'' relies on the physical premise that a single causal event cannot discriminate infinitely many alternatives. This is well-motivated by the operational meaning of ``event,'' but is not a theorem.
    \item \textbf{Elementary $\neq$ composite.} The classical compositeness argument does not cover irreducibly entangled quantum events; Zilly's theorem fills that gap.
\end{enumerate}
\textbf{The upshot:} Readers who accept the structural argument obtain 1-bit quantization as a consequence of Definition~C. Readers who do not may take ``causal events are 1-bit events'' as a well-motivated postulate. The rest of the paper---the $q$-field dynamics (§III--VIII)---is unchanged either way.

\textbf{Finite $c$ + spacing $a$ $\Rightarrow$ locality.}
Causal influence propagates at finite speed $c$ between cells separated by a fundamental spacing $a$. The combination of finite speed and discrete spacing entails that at the fundamental scale, only nearest-neighbor interactions are physically meaningful: no causal event can skip intermediate cells without violating the speed limit $c$. Locality---the principle that direct causal influence occurs only between adjacent cells---is therefore a derived property of the cell graph, not an independent assumption. (The graph structure itself---a regular 3D lattice with uniform spacing---remains a Tier~2 modeling choice; what is derived is the restriction to nearest-neighbor edges.)

\textbf{Undeterminedness between events $\Rightarrow$ superposition.}
Between causal events, a cell's state is not determined by any causal event---an objective undeterminedness. This undeterminedness is the physical substrate of quantum superposition: The quantum reconstruction literature~\cite{Zeilinger1999,Dakic2011,Chiribella2011,Masanes2011,Shaghaghi2023,Zilly:2026} shows that under finite capacity constraints, such undeterminedness takes the mathematical form of complex amplitudes, the Born rule, and unitary dynamics. Superposition is not an additional postulate; it is what the world looks like in the intervals between causal events, and its rigorous mathematical form is supplied by Zilly's theorem. When a causal event occurs, the superposition resolves into a determinate outcome. The full mathematical structure of superposition---complex amplitudes, the Born rule, unitary dynamics---is provided by Zilly's (2026) theorem (Sec.~\ref{sec:foundational}), which derives standard quantum mechanics from finite information capacity and self-referential consistency without assuming superposition. Our framework supplies the physical picture (undeterminedness between events); Zilly's theorem supplies the rigorous mathematical form.

\textbf{Division of labor.}
It is useful to state explicitly what is derived here and what relies on the QM reconstruction:
\begin{itemize}
    \item \textbf{From Definition~C alone:} Directionality (time arrow), the 1-bit quantization of causal events, locality, and the physical picture of superposition as inter-event undeterminedness.
    \item \textbf{From the QM reconstruction literature~\cite{Zeilinger1999,Dakic2011,Chiribella2011,Masanes2011,Shaghaghi2023,Zilly:2026}:} Complex Hilbert space structure, the Born rule, unitary dynamics, and the full mathematical apparatus of quantum mechanics.
\end{itemize}
The two chains converge at the 1-bit cell: Definition~C establishes that causal events are 1-bit events; the reconstruction literature establishes that a finite-capacity system whose states are distinguished by 1-bit events necessarily obeys quantum mechanics. The bridge between the classical causal graph and the quantum Hilbert space is this body of peer-reviewed work, not a single theorem.

In summary, from the single Definition~C we obtain:
\begin{enumerate}
    \item \textbf{Directionality} (asymmetry $\Rightarrow$ time arrow);
    \item \textbf{Capacity bound} (determinate change $\Rightarrow$ discrete $\Rightarrow$ 2 states $\Rightarrow$ 1 bit);
    \item \textbf{Locality} (finite $c$ + spacing $a$ $\Rightarrow$ nearest-neighbor interaction);
    \item \textbf{Superposition} (undeterminedness between causal events $\Rightarrow$ quantum coherence, with mathematical form from the QM reconstruction literature).
\end{enumerate}

\subsection{Foundational Result}
\label{sec:foundational}

We take as given the following result, established independently by the quantum reconstruction literature~\cite{Zeilinger1999,Dakic2011,Chiribella2011,Masanes2011,Shaghaghi2023,Zilly:2026}:

\begin{quote}
\textbf{Finite-capacity quantum reconstruction.} Let a physical system have finite information capacity (one bit per elementary system~\cite{Zeilinger1999,Dakic2011}) and satisfy basic structural constraints (reversibility, locality, purification, or self-referential consistency~\cite{Chiribella2011,Masanes2011,Shaghaghi2023,Zilly:2026}). Then: (1)~the state representation necessarily involves complex numbers; (2)~the probability rule is the Born rule $P(a)=|\langle a|\psi\rangle|^2$; (3)~the dynamics are unitary; and (4)~the resulting theory is standard quantum mechanics.
\end{quote}

We do not re-derive or question this body of work. We accept it as the quantum-mechanical foundation upon which we build the classical gravitational limit.

The architecture of our argument does not depend on any single reconstruction---the classical gravitational limit derived in Sections~\ref{sec:static}--\ref{sec:selforg} depends only on Definition~C and the Tier~2 structural choices, not on the specific mechanism by which quantum mechanics emerges. The key insight we exploit, common to all these reconstructions, is that finite capacity applies at \emph{every} scale---capacity is bounded per cell, not just globally. $\hbar$ enters our framework as the scale of the capacity bound (via phase space quantization), not as a separate axiom.

\subsection{Cell Graph}

Consider a set of $N$ discrete cells arranged on a graph $G=(V,E)$.
Each vertex $i\in V$ is a cell---the site of causal events.
Edges $(i,j)\in E$ are potential causal channels.
The graph has the following properties:

\begin{enumerate}
    \item \textbf{Geometric embedding.} Cells are embedded in $d=3$ spatial dimensions. We take $d=3$ as an empirical input; the theory does not predict spatial dimensionality.
    \item \textbf{Spacing.} Adjacent cells are separated by a fundamental length scale $a$. We identify $a$ with the Planck length $\ell_P\equiv\sqrt{\hbar G/c^3}$. This identification follows from the naturalness requirement that the theory contain no new scales beyond those already present in $c$, $\hbar$, and $G$ (see Sec.~\ref{sec:gamma} for detailed justification).
    \item \textbf{Locality.} Direct causal influence occurs only between adjacent cells (nearest neighbors on the graph). This is a derived consequence of Definition~C: finite propagation speed $c$ and cell spacing $a$ together imply that information cannot skip intermediate cells without violating the speed limit (Sec.~\ref{sec:causal_def}).
    \item \textbf{Speed limit.} Information propagates at most at speed $c$ between adjacent cells, so the minimum time step is $\Delta t = a/c$.
\end{enumerate}

\subsection{Tiered Classification of Assumptions}
\label{sec:tiered}

All inputs to the theory are classified into four tiers, separating irreducible physical principles from modeling choices.

\begin{table}[h]
\centering
\caption{Tiered classification of all assumptions entering the $q$-field framework. Tiers 0--1 are irreducible inputs; Tier~2 choices are protected by universality; Tier~3 quantities are fixed by Tiers~0--2.}
\label{tab:tiers}
\begin{tabular}{p{0.12\columnwidth}p{0.40\columnwidth}p{0.40\columnwidth}}
\toprule
\textbf{Tier} & \textbf{Content} & \textbf{Status} \\
\midrule
Tier 0 & C: A causal event is an asymmetric influence & Irreducible physical definition. \\
(Physical Definition) & producing a determinate change. & \\
\midrule
Tier 1 & $c$: maximum signal speed. & Empirically measured. \\
(Empirical Inputs) & $\hbar$: quantum of action (via QM reconstruction). & Empirically measured. \\
& $G$: Newton's constant. & Empirically measured. \\
& $d=3$: spatial dimensionality. & Empirically observed. \\
\midrule
Tier 2 & Regular 3D lattice. & Simplest choice; universality protects \\
(Modeling Choices) & Nearest-neighbor edges. & macroscopic predictions. \\
& $\phi(q)=-\ln q \in \Phi$. & Any $\phi\in\Phi$ gives same qualitative \\
& & physics (Sec.~\ref{sec:static}). \\
\midrule
Tier 3 & $\gamma_0 = c/a = c/\ell_P$. & Fixed by Tiers 0--2 + naturalness. \\
(Derived Quantities) & $D_0 = c^2/\gamma_0 = c\ell_P$. & Derived from $\gamma_0$. \\
\bottomrule
\end{tabular}
\end{table}

\textbf{Tier 0--1 are irreducible inputs.} They are the physical facts that any theory must accommodate.
\textbf{Tier 2 choices are protected by universality.} Any potential $\phi\in\Phi=\{\phi(q)\mid\phi'(q)<0,\;\phi(q)\to\infty\;\text{as}\;q\to0\}$ gives the same qualitative physics: the static equation is $\nabla^2\phi(q)=0$, the $1/r$ potential emerges, and the resistance to information inflow diverges as $q\to0$.
The specific choice $\phi(q)=-\ln q$ is the simplest member of this universality class.
\textbf{Tier 3 quantities are fixed by Tiers 0--2.} The damping scale $\gamma_0=c/a$ is the only dimensional parameter constructed from the cell spacing $a$ and the speed limit $c$.
Identifying $a$ with $\ell_P$ is a naturalness argument (Sec.~\ref{sec:gamma}): the theory introduces $a$ as the single new scale, and the only natural candidate consistent with the requirement that no new scales appear beyond $c$, $\hbar$, $G$ is the Planck length.

\subsection{The Free-Capacity Fraction $q$}

Each cell $i$ has a total capacity of 1 bit (Definition~C).
At any time $t$, some fraction of this capacity is occupied (storing information) and some is free (available for new information).
Define:

\begin{equation}
q_i(t) \equiv \frac{\text{free capacity of cell } i}{\text{total capacity of cell } i} \in [0, 1].
\label{eq:qdef}
\end{equation}

The free-capacity fraction $q_i$ is the fundamental dynamical variable of our theory.
It is a real scalar field on the discrete graph.
We will show that in the continuum limit, $q$ obeys equations that reproduce both quantum wave dynamics and classical gravitational field equations, depending on the regime.

\subsection{Physical Interpretation}

\begin{itemize}
    \item $q_i = 1$: cell $i$ is completely empty. All capacity is available. This corresponds to asymptotically flat spacetime far from any mass.
    \item $q_i = 0$: cell $i$ is completely full. No capacity remains. This corresponds to the limiting case of infinite information density---the analog of a gravitational singularity, though we will see that $q=0$ is approached asymptotically, never reached at finite $r$.
    \item $0 < q_i < 1$: intermediate occupation. Gradients in $q$ drive information currents, which manifest as gravitational effects.
\end{itemize}

A crucial point: $q$ is \emph{not} a field in spacetime.
Rather, $q$ is a property of the information substrate from which spacetime itself emerges.
The metric structure we perceive as gravity is a coarse-grained description of $q$-field gradients.

\subsection{Coarse-Graining and the Continuum Limit}

Here is a subtlety.
Definition~C entails that each individual cell's causal influence is exactly 1~bit ($s_i\in\{0,1\}$), so a single cell's free-capacity fraction can only be 0 or 1.
The field $q_i(t)\in[0,1]$ used throughout this paper is defined on \emph{mesoscopic blocks} containing $N_b\gg 1$ Planck-scale cells, over which the binary occupancy is averaged:
$q_i \equiv N_b^{-1}\sum_{k\in\text{block}(i)}(1-s_k)$.
For $N_b$ sufficiently large, $q_i$ is effectively continuous, justifying the use of PDEs.
The discrete Bernoulli noise from individual binary cells is suppressed by $1/\sqrt{N_b}$ and is negligible at macroscopic scales.

With $N\sim10^{184}$ Planck cells in the observable universe and $N_b\sim10^{60}$ cells per mesoscopic block (e.g., at the nuclear scale), continuum error estimates are $\mathcal{O}(10^{-123})$ (see SM~B5).
All dynamical equations derived in Sections~\ref{sec:dynamics}--\ref{sec:selforg} describe this coarse-grained field. We assume that coarse-graining commutes with the dynamics to leading order in $a/L$ where $L$ is the macroscopic scale of interest---i.e., that the equations of motion for the block-averaged $q_i$ have the same functional form as those for the underlying binary cells. This assumption is standard in statistical physics (e.g., the Navier--Stokes equations from molecular dynamics) but has not been proved for the $q$-field; a rigorous derivation of the hydrodynamic limit from the discrete cell dynamics is deferred to future work. The binary nature of the 1-bit causal influence manifests only at the Planck scale and does not affect the continuum phenomenology.

%%%%%%%%%%%%%%%%%%%%%%%%%%%%%%%%%%%%%%%%%%%%%%%%%%%%%%%%%%%%%%%%%%%%%%
\section{$q$-Field Dynamics: The Unified Equation}
\label{sec:dynamics}
%%%%%%%%%%%%%%%%%%%%%%%%%%%%%%%%%%%%%%%%%%%%%%%%%%%%%%%%%%%%%%%%%%%%%%

\subsection{Continuity Equation}

Information is conserved in transfer. It moves between cells but is neither created nor destroyed.
For cell $i$, the rate of change of occupied capacity equals the net current into the cell:

\begin{equation}
\frac{\partial}{\partial t}(1 - q_i) = -\sum_{j \in \partial i} J_{ij},
\label{eq:continuity_raw}
\end{equation}

where $\partial i$ denotes the set of cells adjacent to $i$, and $J_{ij}$ is the information current from cell $i$ to cell $j$.
By convention, $J_{ij}>0$ means information flows from $i$ to $j$.
Since $(1-q_i)$ is the occupied fraction, Eq.~\eqref{eq:continuity_raw} simplifies to:

\begin{equation}
\boxed{\frac{\partial q_i}{\partial t} = \sum_{j \in \partial i} J_{ij}}.
\label{eq:continuity}
\end{equation}

The current is antisymmetric: $J_{ij}=-J_{ji}$, reflecting that information leaving $i$ arrives at $j$.
This is the discrete analog of the conservation law $\partial_t q + \nabla\cdot\mathbf{J}=0$ in the continuum.

\subsection{Fundamental Current Equation}

We need a constitutive relation for $J_{ij}$.
The current is driven by gradients in free capacity: information tends to flow from cells with more free capacity to cells with less, as the latter are ``information-rich'' and resist further inflow by the capacity bound.
The simplest local, causal dynamics for the current is a Cattaneo-type relaxation equation~\cite{Cattaneo:1948,Jou:2010}.

\textbf{The nonlinear current equation.}
Four physical constraints must be satisfied:
\begin{enumerate}
    \item The current must vanish when $q_i=q_j$ (no gradient, no flow).
    \item The driving force must diverge as $q_i\to0$ or $q_j\to0$ (no cell can accept information beyond its 1-bit capacity).
    \item The current direction must be downhill in $q$.
    \item \textbf{H-theorem compatibility.} The dynamics must admit a Lyapunov functional $E = \int [\Phi(q) + |\mathbf{J}|^2/(2\kappa)] d^3x$ with $dE/dt \le 0$. By the argument of Sec.~\ref{sec:lyapunov}, this requires $\Phi'(q) = \phi(q)$ where $\phi$ is the potential driving the current. With the Shannon-type information functional $\Phi(q) = q\ln q - q + 1$ (the unique additive form), we obtain $\phi(q) = \ln q$ uniquely.
\end{enumerate}

Constraint~(4) selects the logarithmic gradient from the universality class $\Phi$ as the member compatible with the Shannon information functional $\Phi(q)=q\ln q - q + 1$. Any other $\phi\in\Phi$ would admit an H-theorem with a different $\Phi$ (the antiderivative of $\phi$), but the Shannon form is singled out by the requirement that the Lyapunov functional reduce to the standard information-theoretic entropy in the equilibrium limit. The logarithmic gradient is therefore preferred by the conjunction of H-theorem existence and the Shannon additivity property---not by H-theorem existence alone.
For $q_i,q_j\approx1$, $\ln q_i-\ln q_j\approx q_i-q_j$, recovering the linear approximation---but as $q\to0$, the logarithmic gradient diverges, encoding the capacity bound's resistance to further inflow.

Now the damping.
Occupied cells block information flow, so the effective damping seen by a current $J_{ij}$ should scale with the average occupation of the two cells involved.
We therefore write the damping with $\bar{q}\equiv(q_i+q_j)/2$:

\begin{equation}
\boxed{\frac{\partial J_{ij}}{\partial t} = -\frac{c^2}{a^2}(\ln q_i - \ln q_j) - \gamma_0(1-\bar{q})J_{ij}}.
\label{eq:current_fundamental}
\end{equation}

This is the \textbf{fundamental current equation} of the theory.
One equation, used consistently throughout.
The coupling constant $c^2/a^2$ is the natural frequency scale (inverse light-crossing time squared between adjacent cells).
The damping coefficient $\gamma(q)=\gamma_0(1-q)$ will be derived from the capacity bound in Sec.~\ref{sec:gamma}; for now, $\gamma_0$ is a dimensionful constant with units of inverse time.

Equation~\eqref{eq:current_fundamental} satisfies all three constraints: it vanishes at equal $q$, diverges at $q\to0$, and drives flow downhill.
It is \textbf{not} two separate equations (linear and nonlinear) patched together---it is one unified equation valid in all regimes.

\subsection{Unified Telegraph Equation}

Differentiate the continuity condition with respect to time and substitute the current equation.
The algebra is straightforward but the result is compact:

\begin{align}
\frac{\partial^2 q_i}{\partial t^2} &= \sum_{j \in \partial i} \frac{\partial J_{ij}}{\partial t} \nonumber \\
&= \sum_{j \in \partial i} \left[-\frac{c^2}{a^2}(\ln q_i - \ln q_j) - \gamma_0(1-\bar{q}_{ij})J_{ij}\right] \nonumber \\
&= -\frac{c^2}{a^2} \sum_{j \in \partial i} (\ln q_i - \ln q_j) - \gamma_0 \sum_{j \in \partial i}(1-\bar{q}_{ij})J_{ij}.
\label{eq:telegraph_step}
\end{align}

In the continuum limit $a\to0$, the discrete Laplacian of $\ln q$ becomes:
\begin{equation}
\frac{1}{a^2}\sum_{j \in \partial i}(\ln q_i - \ln q_j) \;\longrightarrow\; -\nabla^2(\ln q).
\end{equation}

For the damping sum, when $q$ varies smoothly across the lattice, $\bar{q}_{ij}\approx q_i$ for all neighbors $j$.
Then:
\begin{equation}
\sum_{j \in \partial i}(1-\bar{q}_{ij})J_{ij} \approx (1-q_i)\sum_{j \in \partial i}J_{ij} = (1-q_i)\frac{\partial q_i}{\partial t}.
\end{equation}

We thus obtain the \textbf{unified telegraph equation}:

\begin{equation}
\boxed{\frac{\partial^2 q}{\partial t^2} + \gamma_0(1-q)\frac{\partial q}{\partial t} = c^2\nabla^2(\ln q)}.
\label{eq:telegraph_unified}
\end{equation}

This is the fundamental dynamical equation of $q$-field theory.
In the derivation we have neglected the cross-term $\sum_j J_{ij}\cdot\nabla(\ln q_i)$, which is $\mathcal{O}(|\nabla q|^2)$. For small perturbations around uniformity this term is subdominant to $(1-q)\partial_t q$; a detailed error bound is given in SM~A3.3.

\subsection{Limiting Cases}

The unified telegraph equation~\eqref{eq:telegraph_unified} interpolates between qualitatively different regimes, all contained within the single equation.
We examine three of them.

\subsubsection{Small-Amplitude Limit: Linear Telegraph}

For perturbations $\delta q$ around a background $q_0$:
\begin{equation}
\ln q = \ln(q_0 + \delta q) \approx \ln q_0 + \frac{\delta q}{q_0}.
\end{equation}

For the important case of vacuum perturbations ($q_0=1$):
\begin{equation}
\ln(1+\delta q) \approx \delta q \quad\Longrightarrow\quad \nabla^2(\ln q) \approx \nabla^2(\delta q).
\end{equation}

The damping term becomes $\gamma_0(1-1-\delta q)\partial_t(\delta q) = -\gamma_0\delta q\,\partial_t\delta q$, which is \textbf{second order} in $\delta q$.
Therefore, to linear order in perturbations around $q=1$:

\begin{equation}
\boxed{\frac{\partial^2 \delta q}{\partial t^2} - c^2\nabla^2\delta q = 0 \quad\Longrightarrow\quad \Box(\delta q) = 0}.
\label{eq:kg_vacuum}
\end{equation}

Perturbations around vacuum propagate as massless waves at speed $c$.
The damping term, being second order in the perturbation, does not appear at linear order.

More generally, for perturbations around $q_0<1$:
\begin{equation}
\frac{\partial^2 \delta q}{\partial t^2} + \gamma_0(1-q_0)\frac{\partial \delta q}{\partial t} - c^2\nabla^2\left(\frac{\delta q}{q_0}\right) = 0,
\label{eq:linear_telegraph_general}
\end{equation}
which is the linear telegraph equation with effective damping $\gamma_{\rm eff}=\gamma_0(1-q_0)$ and an effective diffusion constant $D_{\rm eff}=c^2/[\gamma_0(1-q_0)]$.

\subsubsection{Vacuum Limit: $\gamma\to0$}

In exact vacuum ($q=1$ everywhere), the damping term vanishes identically:
\begin{equation}
\gamma(1) = \gamma_0(1-1) = 0,
\end{equation}
and Eq.~\eqref{eq:telegraph_unified} reduces to:
\begin{equation}
\boxed{\Box(\ln q) = 0},
\label{eq:box_lnq}
\end{equation}
where $\Box\equiv\partial_t^2-c^2\nabla^2$ is the d'Alembertian.
This is the nonlinear wave equation for $\ln q$ in vacuum.
For small perturbations, it reduces to $\Box(\delta q)=0$ as shown above.

\subsubsection{Overdamped Regime: $\gamma\gg\partial_t$}

The damping function $\gamma(q)=\gamma_0(1-q)$ is bounded above by $\gamma_0$ (attained at $q=0$).
There is \textbf{no} ``$\gamma\to\infty$'' limit---$\gamma$ is always finite, with maximum value $\gamma_0=c/\ell_P\approx1.85\times10^{43}$\,s$^{-1}$.

A physically important regime occurs when the inertial term is negligible compared to damping:
\begin{equation}
|\partial_t^2 q| \ll \gamma|\partial_t q|.
\end{equation}
This condition is met on timescales $t\gg\gamma^{-1}$ or, equivalently, for wavenumbers $k\ll k_{\rm cross}\equiv\gamma/(2c)$.
In this \textbf{overdamped regime}, Eq.~\eqref{eq:telegraph_unified} reduces to the effective diffusion equation:

\begin{equation}
\boxed{\frac{\partial q}{\partial t} = D_{\rm eff}\nabla^2(\ln q), \qquad D_{\rm eff} \equiv \frac{c^2}{\gamma_0(1-q)}}.
\label{eq:diffusion}
\end{equation}

The crossover between wave and diffusive behavior is controlled by the dimensionless parameter:
\begin{equation}
R \equiv \frac{\omega}{\gamma} \quad\text{or}\quad \frac{k}{k_{\rm cross}}.
\end{equation}
$R\gg1$: wave regime (inertial term dominates). $R\ll1$: diffusion regime (damping dominates).
This crossover is a sharp prediction of the theory (Sec.~\ref{sec:predictions}).

\textbf{Rigorous analysis.}
The transition from hyperbolic (telegraph) to parabolic (diffusion) behavior is a singular perturbation problem.
As $\epsilon\equiv\gamma^{-1}\to0$, the highest time derivative is multiplied by a small parameter, creating a boundary layer at $t\sim\epsilon$.
For $t\gg\epsilon$, the outer solution obeys the diffusion equation~\eqref{eq:diffusion}; the inner solution matches the initial wavefront.
This structure is standard for telegraph $\to$ diffusion limits~\cite{Jou:2010} and guarantees that information propagates causally (front speed $c$) even in the diffusion-dominated regime.

%%%%%%%%%%%%%%%%%%%%%%%%%%%%%%%%%%%%%%%%%%%%%%%%%%%%%%%%%%%%%%%%%%%%%%
\section{Damping from Capacity: $\gamma(q)=\gamma_0(1-q)$}
\label{sec:gamma}
%%%%%%%%%%%%%%%%%%%%%%%%%%%%%%%%%%%%%%%%%%%%%%%%%%%%%%%%%%%%%%%%%%%%%%

\subsection{Physical Argument}

In Eq.~\eqref{eq:current_fundamental}, the damping term $-\gamma_0(1-\bar{q})J_{ij}$ represents the loss of coherence in the information current.
The physical origin is straightforward: information propagation requires \emph{free capacity} in the receiving cell.
When a cell is occupied (probability $1-q$), it cannot accept new information and therefore blocks the current.

Within any local neighborhood, the fraction of occupied cells encountered by a current is $(1-q)$.
It follows that:

\begin{equation}
\gamma \propto (1-q).
\label{eq:gamma_scaling}
\end{equation}

Damping is not an external friction parameter.
It is an \emph{intrinsic} consequence of the causal definition (Definition~C).
The more occupied the cells, the harder it is for information to flow coherently.
That observation is the central physics of this section.

\subsection{The Linear Ansatz}

The simplest functional form consistent with the physical constraints is:

\begin{equation}
\boxed{\gamma(q) = \gamma_0 (1-q)}.
\label{eq:gamma_linear}
\end{equation}

We require two physical conditions: (i)~$\gamma\to0$ as $q\to1$ (vacuum is undamped---empty cells cannot block), and (ii)~$\gamma>0$ for $q<1$ (occupied cells damp coherent propagation). The simplest monotone ansatz satisfying both constraints is $\gamma_0(1-q)$; any higher-order correction must take the form $(1-q)^n$ with $n\ge2$, which is subdominant when $1-q\ll1$ and does not affect the essential physics. The linear form is thus the leading term in a systematic expansion around the vacuum, not an arbitrary guess.

This satisfies:

\begin{itemize}
    \item \textbf{Vacuum limit:} As $q\to 1$ (completely empty cells), $\gamma\to 0$. In vacuum, there are no occupied cells to block information flow, so damping vanishes entirely. This is a dynamical consequence of $q\to 1$, not a limit we take by hand.
    \item \textbf{Saturation limit:} As $q\to 0$ (completely full cells), $\gamma\to\gamma_0$. All cells are occupied, blocking is maximal, and the current experiences the strongest possible damping.
\end{itemize}

\textbf{Honest statement:} The linear form $\gamma(q)=\gamma_0(1-q)$ is a \emph{modeling ansatz}, not a rigorous derivation from Definition~C.
Nonlinear corrections (e.g., $\gamma\propto(1-q)^\alpha$ with $\alpha\neq 1$) may exist.
The essential point---that $\gamma$ must vanish as $q\to 1$---follows from Definition~C regardless of the specific functional form.
The linear ansatz captures this essential physics with minimal assumptions and is the leading-order term in a $(1-q)$ expansion around vacuum.

\subsection{Naturalness Argument for $\gamma_0$}
\label{sec:naturalness}

$\gamma_0$ has dimensions of inverse time.
The cell spacing $a$ is the \textbf{single new dimensional parameter} introduced by the theory.
Combined with the speed limit $c$, it defines the natural frequency scale:

\begin{equation}
\gamma_0 = \frac{c}{a}.
\label{eq:gamma0_form}
\end{equation}

What fixes $a$?
The theory contains exactly three dimensionful constants from Tier~1: $c$, $\hbar$, and $G$.
If the theory is to contain \textbf{no new scales} beyond those already present in known physics, then $a$ must be constructed from $c$, $\hbar$, and $G$.
Dimensional analysis gives a unique combination with dimensions of length:

\begin{equation}
a = \ell_P \equiv \sqrt{\frac{\hbar G}{c^3}} \approx 1.62 \times 10^{-35}\;\text{m}.
\label{eq:lp_def}
\end{equation}

Consequently:

\begin{equation}
\boxed{\gamma_0 \equiv \frac{c}{\ell_P} = \sqrt{\frac{c^5}{\hbar G}} \approx 1.85 \times 10^{43}\;\text{s}^{-1}}.
\label{eq:gamma0}
\end{equation}

We stress: this is a \textbf{naturalness argument}, not a logical necessity.
The theory introduces $a$ as a free parameter; we \emph{choose} to identify it with $\ell_P$ because doing so introduces no new scales.
This choice has falsifiable consequences: if observations constrain $\gamma_0 \neq c/\ell_P$, the identification is excluded.

In the language of Tier~3 (Table~\ref{tab:tiers}), $\gamma_0=c/\ell_P$ is a derived quantity, fixed by Tier~0--2 inputs plus the naturalness requirement.

\textbf{Dimensionless prefactor.}
Dimensional analysis alone permits an arbitrary $\mathcal{O}(1)$ dimensionless prefactor: $\gamma_0 = \alpha\,c/a$ with $\alpha\sim\mathcal{O}(1)$.
Setting $\alpha=1$ is the simplest choice and is adopted throughout this paper.
All qualitative predictions (Lorentz selection in vacuum, wave--diffusion crossover existence, Newtonian limit recovery, H-theorem) are independent of $\alpha$.
Naturalness prefers $\alpha\sim\mathcal{O}(1)$; the exact value must be determined by observation, making $\alpha$ the single dimensionless free parameter of the theory.

\subsection{Modified Telegraph Equation}

Substituting $\gamma(q)=\gamma_0(1-q)$ into Eq.~\eqref{eq:telegraph_unified}:

\begin{equation}
\boxed{\frac{\partial^2 q}{\partial t^2} + \gamma_0(1-q)\frac{\partial q}{\partial t} = c^2\nabla^2(\ln q)}.
\label{eq:telegraph_gammaq}
\end{equation}

Key properties:
\begin{enumerate}
    \item \textbf{Nonlinearity:} The damping term $\gamma_0(1-q)\partial_t q$ couples the field to its own time derivative. Additionally, the right-hand side $\nabla^2(\ln q)$ provides nonlinear spatial coupling.
    \item \textbf{Vacuum linearization:} For perturbations $\delta q$ around $q=1$, the damping term is $O(\delta q^2)$ and the spatial term linearizes to $\nabla^2\delta q$, giving $\Box(\delta q)=0$ exactly at linear order.
    \item \textbf{Static nonlinearity:} In the static limit $\partial_t\to0$, Eq.~\eqref{eq:telegraph_gammaq} directly yields $\nabla^2(\ln q)=0$, the equation whose solution gives Newtonian gravity (Sec.~\ref{sec:static}).
\end{enumerate}

%%%%%%%%%%%%%%%%%%%%%%%%%%%%%%%%%%%%%%%%%%%%%%%%%%%%%%%%%%%%%%%%%%%%%%
\section{Lorentz Invariance in the Vacuum Limit}
\label{sec:lorentz}
%%%%%%%%%%%%%%%%%%%%%%%%%%%%%%%%%%%%%%%%%%%%%%%%%%%%%%%%%%%%%%%%%%%%%%

\subsection{The Selection Property}

\begin{quote}
\textbf{Vacuum Lorentz Restoration.}
In vacuum, defined as the state $q=1$ (all cells empty), the capacity-dependent damping vanishes: $\gamma(1)=\gamma_0(1-1)=0$.
The dynamics reduce to $\Box(\ln q)=0$ and, for perturbations, $\Box(\delta q)=0$.
Since $\Box$ is invariant under Lorentz transformations, vacuum spacetime is automatically Lorentz invariant.
Lorentz invariance is \textbf{not} imposed as an axiom---it is \textbf{restored} as the ground state of the empty universe.
\end{quote}

We did not \emph{separately} assume Lorentz invariance as an axiom.
The d'Alembertian operator $\Box \equiv \partial_t^2 - c^2\nabla^2$ is built into the fundamental current equation~(\ref{eq:current_fundamental}) via the propagation speed $c$---so Lorentz symmetry with respect to $c$ is a structural feature of the dynamics, not an emergent phenomenon.
The nontrivial claim is narrower: the damping term $\gamma_0(1-q)\partial_t q$, which breaks Lorentz invariance, vanishes identically in vacuum.
Vacuum is therefore the regime where the full Lorentz symmetry of the underlying wave operator is restored, not broken.
The preferred-frame sector ($\gamma>0$) is dynamically confined to matter-filled regions; it is not an independent axiom but a consequence of the same damping that follows from Definition~C.

The mechanism is transparent:
\begin{enumerate}
    \item Definition~C (determinate change implies 1-bit quantization) implies that occupied cells block new causal events.
    \item This blocking produces damping $\gamma\propto(1-q)$.
    \item In vacuum ($q=1$), there are no occupied cells $\Rightarrow$ no blocking $\Rightarrow$ $\gamma=0$.
    \item With $\gamma=0$, the dynamics are $\Box(\ln q)=0$, which is Lorentz invariant.
\end{enumerate}

\textbf{Honest statement:} Lorentz selection here is weaker than the diffeomorphism emergence claimed in some other frameworks.
Our theory always \emph{contained} a Lorentz-invariant sector ($\Box(\ln q)=0$); the result is that vacuum \emph{automatically selects} this sector rather than us discovering it by taking a limit by hand.
The nontrivial claim is that the preferred-frame sector ($\gamma>0$) is dynamically confined to matter-filled regions.

\subsection{Preferred-Frame Effects Are Confined to Matter}

For $q<1$ (anywhere matter is present), $\gamma>0$ and the dynamics are not Lorentz invariant.
The damping term picks out a preferred rest frame---the frame of the information substrate.
This is physically expected: decoherence requires a preferred frame~\cite{Bassi:2013mta}, and here that frame is the rest frame of the cell graph.

Is the resulting Lorentz violation compatible with existing constraints?
For a spherically symmetric mass $M$ at distance $r$, Eq.~\eqref{eq:q_solution} gives:

\begin{equation}
1 - q(r) = 1 - e^{-GM/rc^2} \approx \frac{GM}{rc^2}.
\label{eq:lv_estimate}
\end{equation}

At the solar surface ($M=M_\odot$, $r=R_\odot$):

\begin{equation}
1 - q \approx \frac{GM_\odot}{R_\odot c^2} \approx 2 \times 10^{-6}.
\label{eq:solar_lv}
\end{equation}

So $\gamma/\gamma_0 \approx 2\times 10^{-6}$ at the solar surface.
All existing Lorentz-violation constraints~\cite{Kostelecky:2008,Mewes:2013} are satisfied.
The theory was not tuned to satisfy them---the same $GM/rc^2$ factor that makes gravity weak also makes Lorentz violation small. In the laboratory ($M\sim$ kg, $r\sim$ m), $1-q\sim 10^{-27}$, utterly negligible.

%%%%%%%%%%%%%%%%%%%%%%%%%%%%%%%%%%%%%%%%%%%%%%%%%%%%%%%%%%%%%%%%%%%%%%
\section{Static Limit and Gravitational Potential}
\label{sec:static}
%%%%%%%%%%%%%%%%%%%%%%%%%%%%%%%%%%%%%%%%%%%%%%%%%%%%%%%%%%%%%%%%%%%%%%

Having established that the Lorentz-invariant sector is restored in vacuum and preferred-frame effects are confined to matter, we now turn to the static limit---the regime where the theory makes contact with classical gravity.

\subsection{Unified Static Derivation}

Consider stationary configurations: $\partial_t q=0$ and $\partial_t J_{ij}=0$.
From Eq.~\eqref{eq:current_fundamental} with $\partial_t J_{ij}=0$:

\begin{equation}
0 = -\frac{c^2}{a^2}(\ln q_i - \ln q_j) - \gamma_0(1-\bar{q}_{ij})J_{ij}.
\label{eq:static_J_eq}
\end{equation}

Solving for the equilibrium current:

\begin{equation}
J_{ij} = -\frac{c^2}{\gamma_0(1-\bar{q}_{ij})a^2}(\ln q_i - \ln q_j).
\label{eq:J_static}
\end{equation}

The zero-net-current condition from Eq.~\eqref{eq:continuity} with $\partial_t q_i=0$ gives $\sum_{j\in\partial i}J_{ij}=0$ for all $i$:

\begin{equation}
\sum_{j\in\partial i} \frac{c^2}{\gamma_0(1-\bar{q}_{ij})a^2}(\ln q_i - \ln q_j) = 0.
\label{eq:zero_net_current}
\end{equation}

In the continuum limit $a\to0$, the denominator $(1-\bar{q}_{ij})$ approaches a common value across the neighborhood of $i$ and factors out of the sum to leading order.

\textbf{Honest caveat:} Both derivation paths require the field to be smooth on the scale of the lattice spacing $a$, so that gradient corrections of order $a|\nabla\ln q|$ are negligible. Near mass concentrations where $q$ varies steeply, this approximation receives controlled corrections set by the dimensionless parameter $a|\nabla\ln q|$; the qualitative conclusion ($\nabla^2(\ln q)=0$ in vacuum) is robust.

A cleaner path (detailed in SM, Sec.~A3.5) derives the same result directly from $\partial_t J_{ij}=0$ and $\partial_t q_i=0$, bypassing both the telegraph equation and the parabolic approximation.

Either path yields the same static equation:

\begin{equation}
\sum_{j \in \partial i} (\ln q_i - \ln q_j) = 0.
\label{eq:laplace_ln_disc}
\end{equation}

In the continuum limit $a\to0$:

\begin{equation}
\boxed{\nabla^2(\ln 1/q) = 0}.
\label{eq:laplace_ln}
\end{equation}

This is the fundamental static equation of $q$-field theory, derived directly from the unified dynamics without any patching, separate ``nonlinear correction,'' or ad hoc modification.
The derivation is one continuous argument from Eq.~\eqref{eq:current_fundamental} to Eq.~\eqref{eq:laplace_ln}:

\begin{center}
\small
$\partial_t J_{ij}=0 \;\Rightarrow\; J_{ij} \propto (\ln q_i-\ln q_j)/(1-\bar{q}_{ij}) \;\Rightarrow\; \sum_j J_{ij}=0 \;\Rightarrow\; \sum_j(\ln q_i-\ln q_j)=0 \;\Rightarrow\; \nabla^2(\ln q)=0.$
\end{center}

\textbf{Universality.}
The specific choice $\phi(q)=-\ln q$ belongs to the universality class $\Phi=\{\phi(q) \mid \phi'(q)<0,\;\phi(q)\to\infty\;\text{as}\;q\to0\}$.
Any $\phi\in\Phi$ gives the same qualitative physics: resistance diverges as $q\to0$, the static equation is $\nabla^2\phi(q)=0$, and the $1/r$ potential emerges.
The choice $-\ln q$ is the simplest member of this class and is protected by universality---the macroscopic physics does not depend on microscopic details of $\phi$.

\subsection{Spherically Symmetric Solution}

For spherical symmetry, $\ln(1/q)$ depends only on $r$.
Setting $\nabla^2(\ln 1/q)=0$:

\begin{equation}
\frac{1}{r^2}\frac{d}{dr}\left(r^2 \frac{d}{dr}\ln\frac{1}{q}\right) = 0,
\label{eq:ode}
\end{equation}

which integrates to $r^2\frac{d}{dr}\ln(1/q)=\text{const}$.
Writing the integration constant as $A = GM/c^2$, where $M$ is the source mass and $G$ is Newton's constant---a calibration, not a derivation: the $1/r$ functional form is derived from $\nabla^2(\ln 1/q)=0$ and $d=3$, but the coefficient $GM/c^2$ is fixed by matching to the observed Newtonian force law:

\begin{equation}
r^2 \frac{d}{dr}\ln\frac{1}{q} = \frac{GM}{c^2}.
\label{eq:first_int}
\end{equation}

Integrating again with the boundary condition $q\to1$ as $r\to\infty$:

\begin{equation}
\boxed{q(r) = \exp\left(-\frac{GM}{rc^2}\right)}.
\label{eq:q_solution}
\end{equation}

This is the central result: the free-capacity fraction around a spherically symmetric mass $M$.

\subsection{Recovery of Newtonian Gravity}

For $r\gg GM/c^2$ (the weak-field regime):

\begin{equation}
q(r) = 1 - \frac{GM}{rc^2} + \mathcal{O}\left(\frac{G^2 M^2}{r^2 c^4}\right).
\label{eq:q_expand}
\end{equation}

The connection to the Newtonian potential follows from the weak-field expansion. For $r\gg GM/c^2$, $q(r)\approx 1-GM/rc^2$, so $1-q \approx GM/rc^2$. The gravitational potential is identified as $\Phi(r) \equiv -c^2(1-q(r))$, motivated by the requirement that test particles follow geodesics of the emergent metric (Sec.~\ref{sec:metric}) and that the Newtonian limit $\Phi = -GM/r$ be recovered. In the weak-field limit:

\begin{equation}
\boxed{\Phi(r) = -\frac{GM}{r} + \mathcal{O}(r^{-2})}.
\label{eq:newton}
\end{equation}

Newton's law is recovered in the weak-field, static limit.
We now show that the same information-theoretic principles fix the full spacetime metric, not just the Newtonian potential.

\subsection{Emergent Spacetime via Causal Locking}
\label{sec:metric}

\textbf{Causal locking and $g_{00}$.}
We adopt the \textbf{causal locking hypothesis}: the rate at which proper time advances locally is set by the free-capacity fraction, $d\tau/dt = q$. The physical motivation is that time advances through the occurrence of causal events; in a region where cells are nearly saturated ($q \ll 1$), few causal events can occur per coordinate time step, so proper time advances slowly. This hypothesis is a structural postulate, not a derived result; it is the minimal bridge from the $q$-field to metric structure.
In the rest frame of a static observer, $d\tau^2 = -g_{00}dt^2$, giving $g_{00} = -q^2$.

\textbf{Covariance consistency and $g_{rr}$.}
The spatial metric is fixed by a mathematical consistency requirement rather than a heuristic signal-propagation argument. In isotropic coordinates $ds^2 = -A(r)c^2dt^2 + B(r)(dr^2 + r^2d\Omega^2)$, the static $q$-field equation $\nabla^2(\ln q)=0$ must be the static specialization of the covariant wave equation $\Box_g(\ln q)=0$ in the effective metric. The d'Alembertian in these coordinates gives $\Box_g\phi = \frac{1}{\sqrt{A}B^{3/2}r^2}\partial_r(\sqrt{AB}\,r^2\partial_r\phi)$ for static $\phi(r)$. This reduces to the flat-space Laplacian $\frac{1}{r^2}\partial_r(r^2\partial_r\phi)$ if and only if $\sqrt{AB}=1$, i.e., $B=1/A$. With $A = q^2$ (causal locking), this forces $B = 1/q^2$, yielding $g_{rr} = 1/q^2$ in isotropic coordinates.

This derivation replaces the heuristic mean-hop-time argument with a precise mathematical condition: the static limit of the $q$-field equation is $\nabla^2(\ln q)=0$, which must be the static specialization of the covariant equation in the effective metric. The condition $AB=1$ is the unique choice (in isotropic coordinates) for this to hold, independent of any probabilistic reasoning about cell occupancy.

\textbf{Full static metric.}
With $q(r) = e^{-GM/rc^2}$, the emergent line element is:
\begin{equation}
\boxed{ds^2 = -e^{-2GM/rc^2}c^2dt^2 + e^{2GM/rc^2}(dr^2 + r^2d\Omega^2)}.
\label{eq:metric}
\end{equation}

\textbf{PPN parameters and consistency with GR.}
Expanding to first post-Newtonian order ($U \equiv GM/rc^2 \ll 1$):
$g_{00} = -(1 - 2U + 2U^2 + \ldots)$, $g_{rr} = 1 + 2U + 2U^2 + \ldots$.
In the standard PPN gauge, this yields $\beta = 1$ and $\gamma = 1$---identical to GR at leading order.
The perihelion precession, light deflection, and Shapiro time delay all match GR predictions exactly at this order.
Second-order deviations ($\Delta g_{rr} \sim U^2/2$) are $\sim10^{-16}$ for Mercury, eight orders of magnitude below current MESSENGER precision, but may be accessible to future strong-field observations.

\textbf{Honest statement:} We recover the Newtonian limit of GR, \emph{not} the full Einstein equations.
The $1/r$ potential was not put in by hand---it follows from spherical symmetry, the static equation $\nabla^2(\ln 1/q)=0$, and the boundary condition $q\to1$ at infinity.
The constant $GM/c^2$ is an integration constant; we identify $G$ as Newton's constant and $M$ as the source mass.
The $1/r$ dependence and exponential form are derived, not assumed.

\subsection{The Role of $d=3$: Marginal Dimensionality}
\label{sec:d3}

Spatial dimensionality $d=3$ is taken as an empirical input (Tier~1).
But it plays a distinguished role.
In $d$ dimensions, the static equation $\nabla^2(\ln 1/q)=0$ integrates to $\ln(1/q)\propto r^{-(d-2)}$ (for $d\neq2$) or $\ln(1/q)\propto\ln r$ (for $d=2$).
The gravitational potential is then $\Phi\sim r^{-(d-2)}$.

Only $d=3$ yields the $1/r$ force law that permits stable orbits, long-range structure formation, and the scaling $\Phi\propto M/r$ that we observe.
In this sense, $d=3$ is a \textbf{marginal} dimension for long-range forces: $d<3$ gives forces that fall off too slowly for stable bound systems; $d>3$ gives forces that fall off too rapidly for hierarchical structure formation.
The fact that our universe selects $d=3$ is an empirical fact, but the theory highlights why this particular dimension enables the rich gravitational phenomenology we observe.

\textbf{Matter coupling.}
The static equation $\nabla^2(\ln 1/q)=0$ describes vacuum configurations; the spherically symmetric solution $q(r)=e^{-GM/rc^2}$ is obtained with an integration constant identified as $GM/c^2$ to match the Newtonian limit.
A complete theory must derive how matter (occupied information capacity) sources the $q$-field in general configurations, yielding a Poisson-like equation with a source term $\propto(1-q)/a^3$.
The derivation of the correct dimensionful prefactor and the full coupling to the stress-energy tensor is left to future work; the vacuum solution calibrated to the Newtonian limit suffices for the static, weak-field results derived here.
A natural candidate source term, consistent with the identification of occupied capacity as information charge, is $\nabla^2(\ln 1/q) = (4\pi G/c^2)\,\rho_m$ with $\rho_m \propto (1-q)/a^3$, which reduces to the standard Poisson equation in the weak-field limit.
Establishing this coupling from the microscopic current dynamics is a priority for extending the framework beyond the vacuum Newtonian regime.

%%%%%%%%%%%%%%%%%%%%%%%%%%%%%%%%%%%%%%%%%%%%%%%%%%%%%%%%%%%%%%%%%%%%%%
\section{Self-Organization: Unifying Decoherence and Collapse}
\label{sec:selforg}
%%%%%%%%%%%%%%%%%%%%%%%%%%%%%%%%%%%%%%%%%%%%%%%%%%%%%%%%%%%%%%%%%%%%%%

\subsection{Decomposition of the Static Equation}

Expand the static nonlinear equation~\eqref{eq:laplace_ln}:

\begin{equation}
\nabla^2(\ln q) = \frac{\nabla^2 q}{q} - \frac{|\nabla q|^2}{q^2} = 0.
\label{eq:expand_ln}
\end{equation}

This gives:

\begin{equation}
\boxed{\frac{\nabla^2 q}{q} = \frac{|\nabla q|^2}{q^2}}.
\label{eq:decomp}
\end{equation}

Equation~\eqref{eq:decomp} expresses a balance between two competing processes.
On one side, the diffusion term $\nabla^2 q/q$---the standard Laplacian diffusion that smooths gradients and homogenizes $q$.
This drives the system toward uniform $q$, corresponding to decoherence (loss of structured information).
On the other, the self-organization term $|\nabla q|^2/q^2$---always non-negative, amplifying existing gradients.
This is positive feedback driving structure formation.

Both terms come from the \textbf{same equation}.
Not two mechanisms. Two faces of one information-dynamic process.

\subsection{Regime Crossover}

Which term dominates depends on $q$:

\begin{itemize}
    \item \textbf{$q\approx1$ (low information density):} The denominator $q^2$ in the self-organization term is $\approx1$, but $\nabla q$ is small (near-flat capacity). Diffusion dominates, and the system tends toward uniformity. This is the \textbf{quantum regime}: decoherence drives the system toward classicality.

    \item \textbf{$q\approx0$ (high information density):} The self-organization term diverges as $q^{-2}$, overwhelming diffusion. The system amplifies any existing gradient, driving further collapse. This is \textbf{gravitational collapse}: structure formation as runaway information concentration.
\end{itemize}

In the intermediate regime ($0<q<1$), the balance determines stability.
From quantum decoherence at $q\approx1$ through astrophysical accretion to gravitational collapse at $q\to0$---one unified picture of structure formation across all scales.

\subsection{H-Theorem: Lyapunov Functional for the Nonlinear Theory}

\label{sec:lyapunov}

We prove that the full nonlinear dynamics possess an H-theorem, establishing thermodynamic consistency.
Define the generalized energy functional:

\begin{equation}
E[q,\mathbf{J}] = \int d^3x\,\left[\Phi(q) + \frac{1}{2\kappa}|\mathbf{J}|^2\right],
\label{eq:lyapunov_E}
\end{equation}

where $\kappa\equiv c^2/a^2$ and $\Phi(q)$ satisfies $\Phi'(q)=\ln q$. (The continuum field $\mathbf{J}$ used here absorbs a factor of $a$ relative to the discrete edge currents; see SM for the equivalent formulation in terms of the bare discrete currents.)
Explicitly, $\Phi(q)=q\ln q - q + 1$ (the constant is chosen so that $\Phi(1)=0$, the vacuum energy).

The continuum limit of the fundamental current equation~\eqref{eq:current_fundamental} is:

\begin{equation}
\frac{\partial \mathbf{J}}{\partial t} = -\kappa\nabla(\ln q) - \gamma_0(1-q)\mathbf{J}.
\label{eq:J_continuum}
\end{equation}
(The continuum field $\mathbf{J}$ in this section is oriented opposite to the discrete edge currents $J_{ij}$ of Sec.~\ref{sec:dynamics}, a convention chosen to simplify the Lyapunov functional. The H-theorem is independent of this orientation choice.)

Together with the continuity equation $\partial_t q = -\nabla\cdot\mathbf{J}$, we compute the time derivative of $E$:

\begin{align}
\frac{dE}{dt} &= \int d^3x\,\left[\Phi'(q)\frac{\partial q}{\partial t} + \frac{1}{\kappa}\mathbf{J}\cdot\frac{\partial\mathbf{J}}{\partial t}\right] \nonumber \\
&= \int d^3x\,\Bigl[\ln q\,(-\nabla\cdot\mathbf{J}) \nonumber \\
&\qquad + \frac{1}{\kappa}\mathbf{J}\cdot\bigl(-\kappa\nabla(\ln q) - \gamma_0(1-q)\mathbf{J}\bigr)\Bigr] \nonumber \\
&= \int d^3x\,\Bigl[-\ln q\,\nabla\cdot\mathbf{J} - \mathbf{J}\cdot\nabla(\ln q) - \frac{\gamma_0(1-q)}{\kappa}|\mathbf{J}|^2\Bigr].
\label{eq:dEdt_step1}
\end{align}

Integrating the first two terms by parts:

\begin{equation}
\int d^3x\,\bigl[-\ln q\,\nabla\cdot\mathbf{J} - \mathbf{J}\cdot\nabla(\ln q)\bigr] = -\oint_{\partial V} (\ln q)\,\mathbf{J}\cdot d\mathbf{S}.
\label{eq:boundary_term}
\end{equation}

For isolated systems ($q\to1$, $\mathbf{J}\to0$ at spatial infinity) or periodic boundary conditions, the surface term vanishes.
Then:

\begin{equation}
\boxed{\frac{dE}{dt} = -\int d^3x\,\frac{\gamma_0(1-q)}{\kappa}\,|\mathbf{J}|^2 \;\leq\; 0}.
\label{eq:dEdt_final}
\end{equation}

Since $\gamma_0(1-q)\geq0$, $\kappa>0$, and $|\mathbf{J}|^2\geq0$, we have $dE/dt\leq0$ for all configurations.
Equality holds only when $\mathbf{J}=0$ everywhere (static equilibrium) or $q=1$ everywhere (vacuum).

\textbf{This proves the H-theorem for the nonlinear theory, conditional on two assumptions:} (i)~the specific choice $\Phi(q)=q\ln q - q + 1$ whose derivative matches the logarithmic gradient, and (ii)~the validity of the continuum limit for the configurations considered (the discrete-to-continuum mapping is assumed; a rigorous convergence proof is deferred to future work). Subject to these conditions, the dynamics are irreversible and thermodynamically consistent.
The Lyapunov functional identifies the static configurations ($\mathbf{J}=0$, $\nabla^2(\ln q)=0$) as the attractors of the dynamics, with the exponential solution~\eqref{eq:q_solution} as the minimum free-energy configuration for a given total information charge (mass).

The physical picture is clean.
$\int\Phi(q)\,d^3x$ is the configuration entropy of the $q$-field.
$\int|\mathbf{J}|^2/(2\kappa)\,d^3x$ is the kinetic energy of information currents.
Damping $\gamma_0(1-q)$ dissipates kinetic energy at a rate proportional to local occupation, driving the system toward the entropy-maximizing static configuration.

\subsection{Time-Dependent Stability Analysis}
\label{sec:stability}

The self-organization decomposition~\eqref{eq:decomp} describes the static balance.
For time-dependent configurations, the full telegraph equation governs the competition between diffusion and collapse.
Consider a localized perturbation with characteristic scale $L$ and amplitude $\delta q$.
Three timescales compete:

\begin{enumerate}
    \item \textbf{Wave timescale:} $\tau_{\rm wave} \sim L/c$ --- the light-crossing time.
    \item \textbf{Damping timescale:} $\tau_{\rm damp} \sim 1/[\gamma_0(1-q_0)]$ --- the time for damping to dissipate coherent oscillations.
    \item \textbf{Diffusion timescale:} $\tau_{\rm diff} \sim L^2/D_{\rm eff} = L^2\gamma_0(1-q_0)/c^2$ --- the time for diffusion to smooth gradients.
\end{enumerate}

The dimensionless ratio controlling the dynamics is:

\begin{equation}
R \equiv \frac{\tau_{\rm damp}}{\tau_{\rm wave}} = \frac{c}{\gamma_0(1-q_0)L}.
\label{eq:stability_R}
\end{equation}

\begin{itemize}
    \item $R\gg1$ (small $L$, high $q_0$): wave propagation dominates. Perturbations oscillate and disperse.
    \item $R\ll1$ (large $L$, low $q_0$): diffusion dominates. Perturbations smooth out on timescale $\tau_{\rm diff}$.
    \item The self-organization term $|\nabla q|^2/q^2$ becomes competitive with diffusion when $|\nabla q|/q \gtrsim 1/L$, i.e., when the logarithmic gradient exceeds the inverse scale. This occurs near mass concentrations where $q\ll1$.
\end{itemize}

A perturbation is unconditionally stable against gravitational collapse when $R\gg1$---the wave dynamics disperse the perturbation before self-organization can amplify it.
Conversely, when $R\ll1$ \emph{and} $q\ll1$, the self-organization term dominates both wave dispersion and diffusion, driving runaway collapse.
This provides a quantitative criterion for the onset of gravitational instability in the $q$-field framework.

\subsection{Variational Interpretation}

Equation~\eqref{eq:laplace_ln} can be understood variationally.
Define the information-theoretic action:

\begin{equation}
\mathcal{S}[q] = \int d^3x \, \frac{|\nabla q|^2}{q^2}.
\label{eq:action}
\end{equation}

The Euler--Lagrange equation $\delta\mathcal{S}/\delta q=0$ yields $\nabla^2(\ln q)=0$.
The exponential solution~\eqref{eq:q_solution} extremizes the integrated squared logarithmic gradient.
Physically: nature settles into the configuration that minimizes information gradients for a given total information charge (mass).
The $1/r$ potential is not a force law but a \textbf{minimum free-energy configuration of the $q$-field}, as proven by the H-theorem of Sec.~\ref{sec:lyapunov}.

\subsection{Quantitative vs.\ Qualitative Predictions}
\label{sec:qual_vs_quant}

We must distinguish predictions that depend on the specific choice $\phi(q)=-\ln q$ from those protected by universality.

\paragraph{Qualitative predictions (universal across $\Phi$):}
\begin{itemize}
    \item Damping vanishes in vacuum: $\gamma(q)\to0$ as $q\to1$.
    \item Lorentz invariance is restored in vacuum.
    \item The static equation is $\nabla^2\phi(q)=0$; the $1/r$ potential emerges for any $\phi\in\Phi$.
    \item Diffusion and self-organization are two regimes of a single equation.
    \item The wave--diffusion crossover exists at $k=\gamma/(2c)$.
    \item A soft boundary replaces the hard event horizon.
    \item The H-theorem holds (with a $\phi$-dependent Lyapunov functional).
\end{itemize}

\paragraph{Quantitative predictions (specific to $\phi(q)=-\ln q$):}
\begin{itemize}
    \item The exact static solution: $q(r)=e^{-GM/rc^2}$.
    \item The Newtonian limit: $\Phi(r)=-GM/r + O(r^{-2})$.
    \item The specific form of the matter--$q$-field coupling (source term).
    \item The exact value of $\gamma_0=c/\ell_P$ (depends on the naturalness argument, not $\phi$).
\end{itemize}

Observational tests of quantitative predictions can distinguish $\phi(q)=-\ln q$ from other members of $\Phi$, while qualitative predictions test the entire universality class.

%%%%%%%%%%%%%%%%%%%%%%%%%%%%%%%%%%%%%%%%%%%%%%%%%%%%%%%%%%%%%%%%%%%%%%
\section{Observational Signatures}
\label{sec:predictions}
%%%%%%%%%%%%%%%%%%%%%%%%%%%%%%%%%%%%%%%%%%%%%%%%%%%%%%%%%%%%%%%%%%%%%%

\textbf{Disclaimer:} The distinctive signatures predicted below are not observable with current or planned detectors. The wave--diffusion crossover is at Planckian frequencies; the soft boundary requires next-generation instruments; and GW170817 consistency rests on an unproven conjecture. These signatures are presented as directional targets for future theory development, not as imminent experimental tests.

\subsection{Wave--Diffusion Crossover}

The most distinctive prediction is a frequency-dependent crossover from wave propagation to diffusion in gravitational dynamics.
From the dispersion relation of Eq.~\eqref{eq:telegraph_gammaq} for a plane-wave perturbation $q\propto e^{i(kx-\omega t)}$ on a background $q_0$:

\begin{equation}
\omega^2 + i\gamma_0(1-q_0)\omega = c^2 k^2/q_0,
\label{eq:dispersion}
\end{equation}

the crossover wavenumber is:

\begin{equation}
\boxed{k_{\rm cross} = \frac{\gamma_0(1-q_0)\sqrt{q_0}}{2c}}.
\label{eq:kcross}
\end{equation}

For $k\gg k_{\rm cross}$ (short wavelengths): the $\partial_t^2 q$ term dominates, and $q$ propagates as a damped wave at speed $c$---conventional gravitational wave behavior.
For $k\ll k_{\rm cross}$ (long wavelengths): the $\gamma\partial_t q$ term dominates, and $q$ evolves diffusively---classical gravitational relaxation.

Observable consequences:
\begin{itemize}
    \item \textbf{Frequency-dependent dispersion.} The phase velocity $v_{\rm phase}=c/\sqrt{1+i\gamma/\omega}$ deviates from $c$ at frequencies near and below $\omega_{\rm cross}=\gamma/2$. For a broadband gravitational wave signal (e.g., a binary inspiral), this produces frequency-dependent arrival time delays.

    \item \textbf{Suppression of low-frequency GWs.} Modes with $\omega<\gamma/2$ are overdamped---they decay exponentially without oscillation.
    This suppresses gravitational wave emission below a cutoff frequency, affecting the low-frequency end of the stochastic gravitational wave background.

    \item \textbf{Ringdown deviations.} The quasi-normal mode spectrum of a post-merger compact object is modified: modes above $k_{\rm cross}$ oscillate; modes below decay monotonically.
\end{itemize}

With $\gamma_0=c/\ell_P\sim10^{43}$ s$^{-1}$, the crossover frequency is Planckian and irrelevant for observations.
However, the \emph{effective} damping $\gamma=\gamma_0(1-q)$ is far smaller in realistic environments.
Near a solar-mass object at $r\sim R_\odot$, $\gamma\sim\gamma_0\times2\times10^{-6}\sim4\times10^{37}$ s$^{-1}$, still astronomically large.
For the crossover to enter the observable band (e.g., LISA frequencies $\sim10^{-3}$ Hz), we would need $(1-q)\sim\gamma/\gamma_0\sim10^{-46}$, which occurs only in extremely deep vacuum far from all mass.

This means the \textbf{wave regime} is the generic prediction for all observable gravitational wave phenomena.
The diffusion regime is relevant only in the extreme strong-field regime ($q\ll1$), where it may affect black hole interiors, the very early universe, or the end state of gravitational collapse.

\textbf{Honest statement:} With $\gamma_0$ fixed by the Planck scale, the wave--diffusion crossover is not observable with current or planned detectors.
Its significance is theoretical: it provides a conceptual unification of wave (quantum coherent) and diffusive (classical decoherent) gravitational dynamics within a single equation. Whether this unification is physically meaningful or merely formal is, in the end, a question for experiments we do not yet know how to do.

\subsection{Soft Boundary: No Hard Event Horizon}

In Eq.~\eqref{eq:q_solution}, $q(r)>0$ for all finite $r$:

\begin{equation}
q(r) > 0 \quad \text{for all} \quad r > 0.
\label{eq:soft}
\end{equation}

There is \textbf{no event horizon} in the strict sense.
The free capacity vanishes only asymptotically as $r\to0$.
Matter falling toward $r=0$ encounters an ever-decreasing $q$, approaching but never reaching the capacity bound.
A \textbf{soft boundary}. Gradual saturation of information capacity without sharp causal disconnection.

The Event Horizon Telescope observations of M87*~\cite{EventHorizonTelescope:2019dse} and Sgr A*~\cite{EventHorizonTelescope:2022xnr} reveal dark central shadows consistent with a photon sphere at $\sim2.6\,GM/c^2$ (for M87*).
In the $q$-field picture, the shadow is produced not by an event horizon but by the steep gradient of $q$ near the origin, which traps photons gravitationally.
The profile may differ from the Kerr metric prediction at the percent level.

Distinguishing the soft boundary from the hard Kerr horizon requires:
\begin{itemize}
    \item Higher-resolution EHT observations (next-generation EHT, space VLBI).
    \item Time-domain observations of orbiting hot spots near the ISCO.
    \item Gravitational wave ringdown spectra from binary black hole mergers, whose quasi-normal mode frequencies may differ between a hard horizon and a soft boundary.
\end{itemize}

\subsection{Gravitational Wave Speed and the Milky Way Halo: An Open Consistency Requirement}
\label{sec:gw170817_halo}

In the vacuum limit ($q\to1$, $\gamma\to0$), gravitational waves propagate at exactly $c$.
This is consistent with GW170817/GRB 170817A~\cite{Abbott:2017apx,GW170817grb}, which constrain $-3\times10^{-15}\leq(v_{\rm GW}-c)/c\leq7\times10^{-16}$.

\textbf{Galactic halo propagation: an open theoretical question.}
GW170817 propagated through the Milky Way halo before detection.
In the $q$-field framework, matter-filled regions have $q<1$ and therefore $\gamma>0$.
With $1-q_0 \sim 5\times10^{-7}$ at the Solar position ($r\approx8$\,kpc, $M_{\rm enc}\approx10^{11}M_\odot$), the effective damping is $\gamma_{\rm halo} \approx 9\times10^{36}$\,s$^{-1}$, so $\gamma_{\rm halo}/\omega \sim 10^{34}$ at LIGO frequencies.
In this strongly overdamped regime, freely propagating scalar $q$-field modes ($k < k_{\rm cross}$) have purely imaginary eigenfrequencies---they decay exponentially in time without oscillation.
No coherent scalar signal propagates through the halo.

\textbf{Conjectured resolution.}
If the gravitational wave signal detected by LIGO corresponds to \emph{tensor} perturbations $h_{\mu\nu}$ of an effective metric $g_{\mu\nu}[q]$ that emerges from the static $q$-field background, and if tensor modes decouple from the scalar sector's damping at linear order, then $\Box h_{\mu\nu}=0$ in vacuum and tensor waves propagate at $c$ regardless of the scalar sector's behavior.
This decoupling is a conjecture requiring future proof: the paper has not derived the effective metric or the tensor perturbation equations from the $q$-field.
Until such a derivation exists, the consistency of the $q$-field framework with GW170817 must be regarded as an open theoretical question, not a resolved check.

\subsection{Separation of Qualitative and Quantitative Predictions}

Following the classification of Sec.~\ref{sec:qual_vs_quant}, we separate:

\begin{center}
\begin{tabular}{p{0.45\columnwidth}p{0.45\columnwidth}}
\toprule
\textbf{Qualitative (universal across $\Phi$)} & \textbf{Quantitative ($\phi=-\ln q$ specific)} \\
\midrule
Wave--diffusion crossover exists & Crossover wavenumber $k_{\rm cross}=\gamma/(2c)$ \\
Soft boundary replaces horizon & Exact $q(r)=e^{-GM/rc^2}$ profile \\
Lorentz selected by vacuum & $\gamma_0=c/\ell_P$ (naturalness) \\
$1/r$ potential emerges (any $\phi$) & Newtonian limit $\Phi=-GM/r$ \\
Damping vanishes at $q\to1$ & $v_{\rm GW}/c = 1/\sqrt{q_0}$ (SM D4.2) \\
\bottomrule
\end{tabular}
\end{center}

%%%%%%%%%%%%%%%%%%%%%%%%%%%%%%%%%%%%%%%%%%%%%%%%%%%%%%%%%%%%%%%%%%%%%%
\section{Discussion}
\label{sec:discussion}
%%%%%%%%%%%%%%%%%%%%%%%%%%%%%%%%%%%%%%%%%%%%%%%%%%%%%%%%%%%%%%%%%%%%%%

\subsection{Summary of the Framework}

The framework rests on one definition plus the Tier~1--2 inputs enumerated in Table~\ref{tab:tiers}. Its principal results, in compressed form, are:

\begin{enumerate}
    \item The unified current equation~\eqref{eq:current_fundamental} with logarithmic gradient and $\bar{q}$-dependent damping yields the telegraph equation~\eqref{eq:telegraph_unified}: $\partial_t^2 q + \gamma_0(1-q)\partial_t q = c^2\nabla^2(\ln q)$. This one equation governs all regimes---no patch, no separate linear and nonlinear equations.
    \item Lorentz invariance is restored in vacuum ($q\to1$); preferred-frame effects are confined to matter-filled regions.
    \item The static limit follows directly from $\partial_t J_{ij}=0$: the factor $\gamma_0(1-\bar{q})$ cancels, yielding $\nabla^2(\ln q)=0$, whose spherically symmetric solution $q(r)=e^{-GM/rc^2}$ recovers Newtonian gravity.
    \item The same equation unifies decoherence and gravitational collapse via $\nabla^2 q/q = |\nabla q|^2/q^2$.
    \item An H-theorem is proved (Sec.~\ref{sec:lyapunov}): $dE/dt\leq0$ for the Lyapunov functional $E=\int[\Phi(q)+|\mathbf{J}|^2/(2\kappa)]d^3x$, establishing thermodynamic consistency.
    \item The theory predicts a wave--diffusion crossover and a soft boundary at black holes.
\end{enumerate}

\subsection{The Holographic Tension}
\label{sec:holographic}

Now the hard part.
The 1-bit-per-event bound in Definition~C is a \textbf{local} bound applied to each cell independently.
For a region of size $R$, the total information capacity is proportional to the number of cells in that region, scaling as the volume:

\begin{equation}
S_{\rm max} \sim \left(\frac{R}{a}\right)^3 \propto R^3 \qquad \text{(Volume scaling)}.
\label{eq:volume_scaling}
\end{equation}

This contrasts with the holographic principle~\cite{tHooft:1993,Susskind:1994vu}, which posits that the maximum entropy of a region scales with its surface area:

\begin{equation}
S_{\rm holographic} \sim \frac{A}{4\ell_P^2} \propto R^2 \qquad \text{(Area scaling)}.
\label{eq:area_scaling}
\end{equation}

Volume versus area. This is a qualitative difference in the scaling of maximum information content.

\textbf{Why this matters.}
The holographic principle is supported by two independent lines of evidence:
(1) black hole thermodynamics, where $S_{\rm BH}=A/4\ell_P^2$ is proportional to horizon area, and
(2) the AdS/CFT correspondence, where a gravitational theory in $d+1$ dimensions is equivalent to a nongravitational theory in $d$ dimensions.
If the $q$-field theory is to be taken seriously as a description of gravity, this tension must be resolved.

\textbf{A conjectured resolution.}
A natural pathway out of this tension is that the 1-bit-per-cell bound is an \emph{ultraviolet} capacity limit, while the \emph{infrared} entropy scaling is determined by entanglement structure---not by how many cells exist in a volume, but by how many are pairwise entangled across its boundary.
If each cell within one Planck length of the boundary is maximally entangled with a cell outside, the entanglement entropy scales as the boundary area: $S_{\rm ent}\sim(R/a)^2\ln 2\propto R^2$.
This resolution is a hypothesis requiring formal proof within the $q$-field dynamics; establishing it rigorously is a priority for future work.

\textbf{Honest statement on the tension.}
DGF predicts $S \propto V$ (volume entropy) from the 1-bit-per-cell capacity bound, while the holographic bound requires $S \leq A/4\ell_P^2$ (area entropy).
This tension is real and unresolved at the formal level.
One possible resolution, noted above, is that the holographic bound applies to the entanglement entropy of quantum states, while DGF's entropy counts classical cell occupancy---the two entropies may not be the same quantity.
Whether these two entropies coincide in the quantum limit is an open question tied to O5 (the basis-independent $q$--quantum-coherence mapping).
A second possibility is that the cell graph is not a regular 3D lattice but is dynamically arranged such that the effective information capacity saturates at the area bound; this would require a holographic reformulation of the cell graph.
For the present paper, we flag the tension honestly and treat it as an open problem requiring future resolution.

\subsection{Fair Comparison with HTMAU Using Tiered Methodology}
\label{sec:htmau_comparison}

Bassani and Magueijo (2025)~\cite{Bassani:2025htma} proposed ``How to Make a Universe'' (HTMAU).
In their framework, diffeomorphism invariance emerges from a Markov chain of stochastic choices evolving toward an absorbing state.
Both HTMAU and the present $q$-field theory share a common conviction: spacetime symmetries are emergent rather than fundamental.
The inputs, mechanisms, and outputs differ substantially, however.
To make the comparison systematic, we apply the \textbf{same tiered methodology} (Sec.~\ref{sec:tiered}) to both theories.

\subsubsection{HTMAU's Assumptions in Tiered Classification}

Enumerating HTMAU's specific assumptions as presented in Ref.~\cite{Bassani:2025htma}:

\begin{enumerate}
    \item \textbf{Preferred foliation.} Spacetime has a preferred time foliation---a global time coordinate that singles out a rest frame. This is an irreducible structural input (Tier~0 equivalent).
    \item \textbf{Global/local variable split.} The universe's state is partitioned into ``global'' variables (constants of nature) and ``local'' variables (field configurations on the foliation). This is a structural choice (Tier~2 equivalent).
    \item \textbf{Matter-gravity separation.} The total action separates into a gravitational sector and a matter sector, with the matter sector treated as a perturbation. This is a modeling choice (Tier~2 equivalent).
    \item \textbf{Local diffeomorphism invariance is imposed before the absorbing state.} The Markov chain is constructed to \emph{preserve} local spatial diffeomorphisms at each step, with time reparameterization invariance emerging only in the absorbing state. This means diffeomorphism invariance is partially \emph{assumed}, not fully derived (Tier~1 equivalent---an empirical input about the symmetry structure).
    \item \textbf{Constants as time-conjugate variables.} The ``global'' variables are canonically conjugate to time, with Poisson bracket structure imposed. This is a structural choice (Tier~2 equivalent).
    \item \textbf{Beta potential functions.} Two $\beta$ potential functions ($\beta_{\rm local}$ and $\beta_{\rm global}$) control the mutation rule of the Markov chain. Their specific forms are free functions (Tier~2 structural choices with continuous parameters).
    \item \textbf{Markov chain mutation rule.} The specific mutation rule $X\to X+\{X,H\}_{\rm local}\delta t + \{\beta_{\rm global},X\}\delta W$ is a structural ansatz (Tier~2 equivalent).
\end{enumerate}

\subsubsection{Tiered Comparison}

Table~\ref{tab:compare_htmau} compares the two frameworks using the same tiered methodology.

\begin{table}[h]
\centering
\caption{Tiered comparison of emergent-spacetime frameworks. Both theories are classified using the same four-tier methodology (Sec.~\ref{sec:tiered}).}
\label{tab:compare_htmau}
\begin{tabular}{p{0.22\columnwidth}p{0.35\columnwidth}p{0.35\columnwidth}}
\toprule
& \textbf{HTMAU (Bassani \& Magueijo 2025)} & \textbf{This work ($q$-field)} \\
\midrule
\textbf{Tier 0} & Causality (via Markov ordering) & C: Causal event is an \\
(Physical & & asymmetric influence producing \\
Axioms / & Preferred foliation & a determinate change \\
Definition) & (2 axioms) & (1 definition) \\
\midrule
\textbf{Tier 1} & $c$, $\hbar$, $G$ & $c$, $\hbar$, $G$, $d=3$ \\
(Empirical & Local diffeo invariance & \\
Inputs) & (imposed pre-absorbing state) & \\
& (4 inputs) & (4 inputs) \\
\midrule
\textbf{Tier 2} & Global/local variable split & Regular 3D lattice \\
(Modeling & Matter-gravity separation & Nearest-neighbor edges \\
Choices) & Constants as time-conjugates & $\phi(q)=-\ln q\in\Phi$ \\
& 2 $\beta$ potential functions & \\
& Markov mutation rule & \\
& (5 structural choices) & (3 structural choices) \\
\midrule
\textbf{Tier 3} & Mutation rate parameters & $\gamma_0=c/\ell_P$ \\
(Derived) & $\beta$ function parameters & $D_0=c^2/\gamma_0=c\ell_P$ \\
& (2+ free parameters) & (0 additional parameters) \\
\midrule
\textbf{Emergent} & Diffeomorphism invariance & Lorentz invariance \\
\textbf{symmetry} & (full GR symmetry) & (SR symmetry) \\
\midrule
\textbf{Emergence} & Stochastic: Markov chain & Deterministic: \\
\textbf{mechanism} & $\to$ absorbing state & $q\to1 \Rightarrow \gamma\to0$ \\
\midrule
\textbf{Static} & FRW matter production rate & $\nabla^2(\ln 1/q)=0$ \\
\textbf{equation} & (cosmological) & $q(r)=e^{-GM/rc^2}$ \\
\midrule
\textbf{Testable} & Near-diffeo invariant & Wave--diffusion crossover, \\
\textbf{predictions} & cosmology & soft boundary, $v_{\rm GW}$ correction \\
\bottomrule
\end{tabular}
\end{table}

\subsubsection{Interpretation}

Both frameworks share a philosophy of emergent spacetime.
HTMAU achieves the stronger result---diffeomorphism invariance, the full symmetry of GR---but requires substantially more input: a preferred foliation (an additional Tier~0 axiom), local diffeomorphism invariance imposed pre-emergence (a Tier~1 input), and five Tier~2 structural choices including two free $\beta$ functions.
The $q$-field framework achieves the weaker result---Lorentz invariance, sufficient for special relativity and the Newtonian limit---but with fewer assumptions at every tier.

Is diffeomorphism invariance necessary at the fundamental level, or is it sufficient to show that Lorentz invariance is restored in vacuum and the Newtonian limit is recovered?
If the latter, the economy of assumptions favors the present approach.
If the former, HTMAU's richer structure is necessary and the present theory must be extended---a direction discussed in the open directions below.

We emphasize: this comparison is methodological, not adversarial.
Both approaches contribute to understanding spacetime as emergent.
The tiered methodology provides a systematic framework for comparing economy and explanatory power across different approaches.

\subsection{Honest Assessment of Limitations}

\begin{enumerate}
    \item \textbf{$\gamma(q)=\gamma_0(1-q)$ is a modeling ansatz.} The linear form is the simplest model consistent with Definition~C but is not rigorously derived from it. Nonlinear corrections $\gamma\propto(1-q)^\alpha$ with $\alpha\neq1$ are possible. The essential prediction---$\gamma\to0$ as $q\to1$---is robust regardless of the functional form.

    \item \textbf{The choice $\phi(q)=-\ln q$ is protected by universality but not unique.} Any $\phi\in\Phi$ gives the same qualitative physics. The $-\ln q$ form is the simplest and most natural, but the universality argument must be made rigorous through renormalization-group analysis.

    \item \textbf{We recover only the Newtonian limit of GR.} The full Einstein equations, with their nonlinear self-interaction, frame-dragging, and cosmological constant, are not derived. Extending the $q$-field framework to full GR requires deriving the right-hand side of $\nabla^2(\ln 1/q)=(\text{source term})$ for general stress-energy configurations, and showing that the result matches the Einstein tensor in the appropriate limit.

    \item \textbf{Matter coupling is not fully derived.} We assumed spherical symmetry and identified an integration constant with $GM/c^2$. A complete theory must derive how $T_{\mu\nu}$ sources $q$ in general configurations. The matter--$q$-field coupling proposed in Sec.~\ref{sec:static} is a first step but requires rigorous derivation from the microscopic dynamics.

    \item \textbf{Spatial dimensionality is an input.} $d=3$ is taken as an empirical fact, not a prediction. The $1/r$ potential follows from $d=3$; in $d$ dimensions the static solution would be $q\propto\exp(-\text{const}/r^{d-2})$, yielding a $1/r^{d-2}$ force law.

    \item \textbf{The holographic bound is not satisfied.} As discussed in Sec.~\ref{sec:holographic}, the 1-bit-per-cell bound implies volume scaling $S\sim R^3$, contradicting the holographic area scaling $S\sim R^2$. This is a serious limitation that requires resolution through nonlocal entanglement, holographic cell arrangement, or dynamical saturation constraints.

    \item \textbf{The $q$--quantum mapping is basis-dependent (Open Problem O5).}
    The $q$-field dynamics derived in Secs.~\ref{sec:dynamics}--\ref{sec:selforg} are classical field equations.
    Their connection to quantum mechanics relies on the quantum reconstruction literature~\cite{Zeilinger1999,Dakic2011,Chiribella2011,Masanes2011,Shaghaghi2023,Zilly:2026} as an external input, not on Definition~C alone.
    Within the classical field description, the results are self-contained: the unified telegraph equation, vacuum Lorentz restoration, and Newtonian gravity are derived from the $q$-field equations without quantum interpretation.
    The quantum interpretation of $q$ as a coherence measure requires the cell-basis mapping via the $\ell_1$-norm coherence measure~\cite{Baumgratz2014}; a basis-independent quantum interpretation remains Open Problem~O5.
    Results in Secs.~\ref{sec:dynamics}--\ref{sec:selforg} do not require O5; they concern the classical field $q(\mathbf{x},t)$ and its dynamics.

    \item \textbf{Lorentz selection is weaker than diffeomorphism emergence.} Our theory always contained a Lorentz-invariant sector; the result is that vacuum selects it. HTMAU's diffeomorphism emergence is a stronger result. Whether this matters depends on whether full diffeomorphism invariance is needed at the fundamental level or only in the classical gravitational limit.

    \item \textbf{Tensor vs.\ scalar sector distinction.} The resolution of the GW170817 halo tension (Sec.~\ref{sec:gw170817_halo}) relies on a distinction between scalar $q$-field perturbations (subject to damping) and tensor metric perturbations (propagating at $c$). A complete theory must derive the coupling between the scalar and tensor sectors from first principles, showing that the decoupling at linear order is robust.
\end{enumerate}

\subsection{Open Directions}

\begin{enumerate}
    \item \textbf{Full GR recovery.} The most important open problem is extending the $q$-field framework to recover the full Einstein equations. The nonlinear structure of $\nabla^2(\ln 1/q)=0$ hints at a deeper connection to the Ricci tensor, but the mapping is not yet known. A Poisson-type equation for $\ln q$ sourced by $\propto(1-q)/a^3$ may relate to the $G_{00}$ component of the Einstein tensor in the Newtonian gauge.

    \item \textbf{Holographic resolution.} Resolving the volume-vs.-area scaling tension (Sec.~\ref{sec:holographic}) is essential. This may involve replacing the regular 3D lattice with a holographic cell arrangement, or deriving an effective area-law bound from the dynamics of entangled cells.

    \item \textbf{Cosmology.} Global expansion corresponds to $q\to1$ (increasing free capacity), which is natural if the universe's information content grows more slowly than its capacity. This may connect to dark energy. The $q$-field should have cosmological solutions that can be compared with $\Lambda$CDM.

    \item \textbf{Entanglement and ER = EPR.} Zilly's QM foundation involves entanglement. In the $q$-field picture, two entangled cells share information, modifying their local $q$. The connection between entanglement structure and spacetime geometry~\cite{Maldacena:2013je} may be derivable in this framework.

    \item \textbf{Quantum gravity regime.} When $q\to0$, the capacity bound is nearly saturated. Zilly's finite-capacity QM predicts strong quantum effects. The $q$-field dynamics in this regime may describe quantum gravitational phenomena without a full theory of quantum gravity.

    \item \textbf{Numerical simulations.} Simulating the $q$-field dynamics in binary merger scenarios and comparing waveforms with GR predictions. The soft boundary and nonlinear damping should produce distinguishable ringdown signals.

    \item \textbf{Experimental constraints on $\gamma_0$ and $\phi(q)$.} While we fix $\gamma_0=c/\ell_P$ from naturalness and $\phi(q)=-\ln q$ from simplicity, both choices are falsifiable. If observations constrain $\gamma_0$ to be smaller (e.g., from the absence of anomalous damping in binary pulsar timing), the identification with the Planck scale would need revision. Similarly, deviations from the $e^{-GM/rc^2}$ profile would constrain $\phi(q)$.

    \item \textbf{Rigorous universality proof.} The claim that any $\phi\in\Phi$ gives the same qualitative physics must be made rigorous. This requires a renormalization-group analysis showing that $\phi(q)=-\ln q$ is the infrared fixed point of the universality class, or a constructive proof that macroscopic observables (potential, force law, stability) are independent of the microscopic $\phi$.
\end{enumerate}

\subsection*{A deeper picture}

It is worth sketching a possible deeper picture, whose precise mathematical development lies beyond the scope of this paper. Zilly's starting point---a universe of perfectly distinguished causal events, where every pair carries unit information---describes a maximally symmetric regime of mutual causality. Within this regime, no event is privileged over any other; there is no arrow of time and no distinction between cause and effect. This is the structure of a Euclidean spacetime: all directions equivalent, pure distinguishability without dynamics.

Self-referential consistency, however, demands that the global state also distinguish itself. A perfectly symmetric mutual-causality structure cannot satisfy this requirement---it provides no preferred frame, no distinguished event, no fact of the matter about ``what the universe is like now.'' The act of self-reference breaks the perfect symmetry. A direction is selected. The symmetry-broken ground state that emerges is precisely the $q=1$ vacuum identified in this work: zero occupied capacity, the unique configuration in which no information is stored and all capacity remains available. The first asymmetric causal event---the deposition of one bit into a single cell---then initiates the causal history that the $q$-field dynamics describe. In this sense, the framework developed here may be understood as the effective description of a symmetry-breaking transition from a pre-causal, perfectly symmetric phase to the causal universe we inhabit. We defer the rigorous formulation of this picture to future work.

\subsection*{Closing remarks}

We have shown that a single definition---a causal event is an asymmetric influence producing a determinate change in a discrete cell---together with a small set of structural choices, yields a dynamical equation that interpolates continuously between quantum decoherence and Newtonian gravity. The static limit recovers the $1/r$ functional form of the Newtonian potential; the coefficient $G$ is calibrated, not predicted. The Lorentz-invariant sector is present from the start and restored in vacuum by the vanishing of state-dependent damping. An H-theorem establishes thermodynamic consistency. The decomposition $\nabla^2 q/q = |\nabla q|^2/q^2$ unifies diffusion and self-organization as two faces of one equation.

The theory is honest about what it does not do: it recovers only the Newtonian limit of GR, it leaves the holographic tension unresolved, and its observational signatures lie beyond current experimental reach. Its contribution is not a final theory but a proof of concept: gravity-like and decoherence-like behavior can coexist within a single information-dynamic equation, with a well-defined inventory of what is derived, what is chosen, and what remains open. The nine limitations enumerated in the Abstract are not defects to be hidden but signposts for the work ahead.

%%%%%%%%%%%%%%%%%%%%%%%%%%%%%%%%%%%%%%%%%%%%%%%%%%%%%%%%%%%%%%%%%%%%%%
\section*{Acknowledgments}
%%%%%%%%%%%%%%%%%%%%%%%%%%%%%%%%%%%%%%%%%%%%%%%%%%%%%%%%%%%%%%%%%%%%%%

We thank J.~Zilly for the machine-checked proof in Lean~4 that finite information capacity plus self-referential consistency yields quantum mechanics~\cite{Zilly:2026}, complementing the peer-reviewed reconstruction literature~\cite{Zeilinger1999,Dakic2011,Chiribella2011,Masanes2011,Shaghaghi2023} that provides this work's definitional grounding.
We thank participants of the Emergent Gravity workshop for discussions.
We thank Bassani and Magueijo for making their HTMAU manuscript available, which provided a valuable benchmark against which to measure the economy of our own assumptions and motivated the tiered classification methodology.

%%%%%%%%%%%%%%%%%%%%%%%%%%%%%%%%%%%%%%%%%%%%%%%%%%%%%%%%%%%%%%%%%%%%%%
\bibliographystyle{apsrev4-2}
\begin{thebibliography}{99}

\bibitem{Zilly:2026}
J.~Zilly, ``Existence as Distinguishability: Quantum Mechanics from Finite Graded Equality,''
arXiv:2603.11900 (2026).
[arXiv preprint; independent verification pending. Not re-derived in this work.
See also the peer-reviewed reconstruction literature: Zeilinger (1999), Dakic \& Brukner (2011),
Chiribella, D'Ariano, \& Perinotti (2011), Masanes \& M\"uller (2011), and Shaghaghi (2023),
which independently establish that finite information capacity plus structural constraints yields
standard quantum mechanics.]

\bibitem{Zeilinger1999}
A.~Zeilinger, ``A Foundational Principle for Quantum Mechanics,''
Found.\ Phys.\ \textbf{29}, 631--643 (1999).

\bibitem{Dakic2011}
B.~Daki\'c and \v{C}.~Brukner, ``Quantum Theory and Beyond: Is Entanglement Special?,''
in \emph{Deep Beauty: Understanding the Quantum World through Mathematical Innovation},
ed.\ H.~Halvorson (Cambridge University Press, 2011), pp.\ 365--392.

\bibitem{Chiribella2011}
G.~Chiribella, G.~M.~D'Ariano, and P.~Perinotti,
``Informational derivation of quantum theory,''
Phys.\ Rev.\ A \textbf{84}, 012311 (2011).

\bibitem{Masanes2011}
L.~Masanes and M.~P.~M\"uller,
``A derivation of quantum theory from physical requirements,''
New J.\ Phys.\ \textbf{13}, 063001 (2011).

\bibitem{Shaghaghi2023}
M.~Shaghaghi, ``The Physical Foundation of Quantum Theory,''
Found.\ Phys.\ \textbf{53}, 1--36 (2023).

\bibitem{Bekenstein1973}
J.~D.~Bekenstein, ``Black Holes and Entropy,''
Phys.\ Rev.\ D \textbf{7}, 2333--2346 (1973).

\bibitem{Verlinde2017}
E.~P.~Verlinde, ``Emergent Gravity and the Dark Universe,''
SciPost Phys.\ \textbf{2}, 016 (2017).

\bibitem{Shannon1948}
C.~E.~Shannon, ``A Mathematical Theory of Communication,''
Bell Syst.~Tech.~J.~\textbf{27}, 379--423, 623--656 (1948).

\bibitem{Cover2006}
T.~M.~Cover and J.~A.~Thomas, \emph{Elements of Information Theory}, 2nd~ed.\ (Wiley, 2006).

\bibitem{Neukart:2024qmm}
F.~Neukart, R.~Brasher, and E.~Marx,
``The Quantum Memory Matrix: A Unified Framework for the Black Hole Information Paradox,''
Entropy \textbf{26}, 1039 (2024).
[Independent validation of the Planck-cell paradigm.]

\bibitem{Green:1987sp}
M.~B.~Green, J.~H.~Schwarz, and E.~Witten,
\textit{Superstring Theory},
Cambridge University Press (1987).

\bibitem{Polchinski:1998rq}
J.~Polchinski,
\textit{String Theory},
Cambridge University Press (1998).

\bibitem{Rovelli:2004tv}
C.~Rovelli,
\textit{Quantum Gravity},
Cambridge University Press (2004).

\bibitem{Ashtekar:2004eh}
A.~Ashtekar and J.~Lewandowski,
``Background independent quantum gravity: A status report,''
Class.\ Quant.\ Grav.\ \textbf{21}, R53 (2004).

\bibitem{Ambjorn:2001cv}
J.~Ambjorn, J.~Jurkiewicz, and R.~Loll,
``Dynamically triangulating Lorentzian quantum gravity,''
Nucl.\ Phys.\ B \textbf{610}, 347 (2001).

\bibitem{Weinberg:1979}
S.~Weinberg,
``Ultraviolet divergences in quantum theories of gravitation,''
in \textit{General Relativity: An Einstein Centenary Survey},
Cambridge University Press (1979).

\bibitem{Reuter:2012}
M.~Reuter and F.~Saueressig,
``Quantum Einstein gravity,''
New J.\ Phys.\ \textbf{14}, 055022 (2012).

\bibitem{Verlinde:2010hp}
E.~P.~Verlinde,
``On the origin of gravity and the laws of Newton,''
JHEP \textbf{04}, 029 (2011).

\bibitem{Jacobson:1995ab}
T.~Jacobson,
``Thermodynamics of spacetime: The Einstein equation of state,''
Phys.\ Rev.\ Lett.\ \textbf{75}, 1260 (1995).

\bibitem{Padmanabhan:2010}
T.~Padmanabhan,
``Thermodynamical aspects of gravity: New insights,''
Rep.\ Prog.\ Phys.\ \textbf{73}, 046901 (2010).

\bibitem{Bassani:2025htma}
P.~M.~Bassani and J.~Magueijo,
``How to make a universe,''
Phys.\ Rev.\ D \textbf{111}, 103529 (2025).

\bibitem{Cattaneo:1948}
C.~Cattaneo,
``Sulla conduzione del calore,''
Atti Sem.\ Mat.\ Fis.\ Univ.\ Modena \textbf{3}, 83 (1948).

\bibitem{Jou:2010}
D.~Jou, J.~Casas-V{\'a}zquez, and G.~Lebon,
\textit{Extended Irreversible Thermodynamics},
Springer (2010).

\bibitem{Abbott:2017apx}
B.~P.~Abbott \textit{et al.} (LIGO Scientific and Virgo Collaborations),
``GW170817: Observation of gravitational waves from a binary neutron star inspiral,''
Phys.\ Rev.\ Lett.\ \textbf{119}, 161101 (2017).

\bibitem{GW170817grb}
B.~P.~Abbott \textit{et al.},
``Gravitational waves and gamma-rays from a binary neutron star merger: GW170817 and GRB 170817A,''
Astrophys.\ J.\ Lett.\ \textbf{848}, L13 (2017).

\bibitem{EventHorizonTelescope:2019dse}
Event Horizon Telescope Collaboration,
``First M87 Event Horizon Telescope results. I. The shadow of the supermassive black hole,''
Astrophys.\ J.\ Lett.\ \textbf{875}, L1 (2019).

\bibitem{EventHorizonTelescope:2022xnr}
Event Horizon Telescope Collaboration,
``First Sagittarius A* Event Horizon Telescope results. I. The shadow of the supermassive black hole in the center of the Milky Way,''
Astrophys.\ J.\ Lett.\ \textbf{930}, L12 (2022).

\bibitem{Bassi:2013mta}
A.~Bassi, K.~Lochan, S.~Satin, T.~P.~Singh, and H.~Ulbricht,
``Models of wave-function collapse, underlying theories, and experimental tests,''
Rev.\ Mod.\ Phys.\ \textbf{85}, 471 (2013).

\bibitem{Maldacena:2013je}
J.~Maldacena and L.~Susskind,
``Cool horizons for entangled black holes,''
Fortsch.\ Phys.\ \textbf{61}, 781 (2013).

\bibitem{Will:2014}
C.~M.~Will,
``The confrontation between general relativity and experiment,''
Living Rev.\ Relativ.\ \textbf{17}, 4 (2014).

\bibitem{Kostelecky:2008}
V.~A.~Kosteleck{\'y} and N.~Russell,
``Data tables for Lorentz and CPT violation,''
Rev.\ Mod.\ Phys.\ \textbf{83}, 11 (2011).

\bibitem{Mewes:2013}
M.~Mewes,
``Tests of Lorentz invariance,''
in \textit{Proceedings of the 6th Meeting on CPT and Lorentz Symmetry},
World Scientific (2013).

\bibitem{tHooft:1993}
G.~'t Hooft,
``Dimensional reduction in quantum gravity,''
arXiv:gr-qc/9310026 (1993).

\bibitem{Susskind:1994vu}
L.~Susskind,
``The world as a hologram,''
J.\ Math.\ Phys.\ \textbf{36}, 6377 (1995).

\end{thebibliography}

\end{document}
